% ═══════════════════════════════════════════════════════════════════
%
%   RAPPORT DE PROJET DE FIN D'ANNÉE
%   Yani Concept by Fati — application mobile, API et back-office
%   pour un centre de beauté
%
%   Compilation :  latexmk -pdf main.tex
%              ou :  pdflatex main && pdflatex main && pdflatex main
%                    (trois passes : table des matières, listes, renvois)
%
%   Les figures manquantes s'affichent en cadres d'attente : déposer le
%   fichier dans figures/ puis recompiler suffit à les remplacer.
%   Voir figures/README.md pour la liste complète des visuels attendus.
%
% ═══════════════════════════════════════════════════════════════════
\documentclass[a4paper,12pt,oneside]{report}

% ═══════════════════════════════════════════════════════════════════
%  ▼▼▼  TOUT CE QUI EST À RENSEIGNER TIENT DANS CE BLOC  ▼▼▼
% ═══════════════════════════════════════════════════════════════════
% Le titre est composé dans un cadre de 9,6 cm : il se replie tout seul,
% aucun retour à la ligne à placer à la main.
\newcommand{\nouvelleTitre}{Conception et réalisation d'une application
mobile, d'une API et d'un back-office pour un centre de beauté}

\newcommand{\nouvelleNature}{Projet de Fin d'Année}
\newcommand{\nouvelleFiliere}{Cybersécurité et Infrastructures Réseau}
\newcommand{\nouvelleAnnee}{2025--2026}

% Les noms ci-dessous servent aussi aux remerciements et au tableau des rôles
% SCRUM. Tant qu'ils ne sont pas renseignés, ils s'affichent en gris italique
% — impossible de les oublier à la relecture.
\newcommand{\nouvelleAuteur}{\aCompleter{Prénom NOM}}
\newcommand{\nouvelleEncadrantPedagogique}{\aCompleter{Pr. Prénom NOM}}
\newcommand{\nouvelleEncadrantEntreprise}{\aCompleter{Mme/M. Prénom NOM}}

\newcommand{\nouvelleEtablissement}{École Marocaine des Sciences de l'Ingénieur}
\newcommand{\nouvelleVille}{Marrakech}
\newcommand{\nouvelleTitreCourt}{Yani Concept by Fati}
% ═══════════════════════════════════════════════════════════════════
%  ▲▲▲  FIN DU BLOC À RENSEIGNER  ▲▲▲
% ═══════════════════════════════════════════════════════════════════

% ═══════════════════════════════════════════════════════════════════
%  Préambule — paquets, style, macros
%  Rapport de projet de fin d'études — Yani Concept by Fati
% ═══════════════════════════════════════════════════════════════════

% ── Langue et encodage ───────────────────────────────────────────
\usepackage[utf8]{inputenc}
\usepackage[T1]{fontenc}
\usepackage[french]{babel}
\usepackage{lmodern}
\usepackage{textcomp}

% Babel français met une espace fine avant « : » et transforme les listes.
% On garde ses règles typographiques, mais on lui interdit de redéfinir les
% puces des listes, pour que itemize reste homogène dans tout le document.
\frenchsetup{StandardLists=true}

% ── Mise en page ─────────────────────────────────────────────────
\usepackage[a4paper,top=2.5cm,bottom=2.5cm,left=3cm,right=2.5cm]{geometry}
\usepackage{setspace}
\onehalfspacing
\usepackage{parskip}
\setlength{\parindent}{0pt}

% ── Graphismes et couleurs ───────────────────────────────────────
\usepackage{graphicx}
\usepackage{xcolor}
\usepackage{float}
\usepackage{tikz}

% ── Tableaux ─────────────────────────────────────────────────────
\usepackage{array}
\usepackage{booktabs}
\usepackage{tabularx}
\usepackage{longtable}
\usepackage{multirow}
\usepackage{colortbl}

% ── Divers ───────────────────────────────────────────────────────
\usepackage{amssymb}
\usepackage{enumitem}
\usepackage{fancyhdr}
\usepackage{titlesec}
\usepackage[font=small,labelfont=bf,labelsep=colon]{caption}

% Babel français légende les tableaux « Table ». L'usage académique marocain,
% et le modèle de rapport suivi ici, disent « Tableau ». La redéfinition passe
% par \addto\captionsfrench, sans quoi babel la réécrirait au \begin{document}.
\addto\captionsfrench{%
  \renewcommand{\tablename}{Tableau}%
  \renewcommand{\listtablename}{Liste des tableaux}%
  \renewcommand{\listfigurename}{Liste des figures}%
}
\usepackage{emptypage}
\usepackage{lastpage}
\usepackage{listings}

% hyperref en dernier (il redéfinit beaucoup de commandes)
\usepackage[hidelinks,breaklinks=true]{hyperref}

% ═══════════════════════════════════════════════════════════════════
%  Charte graphique — reprise de la palette réelle de l'application
%  (or #D8A848 sur crème #F9F6F0, cf. mobile/src/theme/colors.ts)
% ═══════════════════════════════════════════════════════════════════
\definecolor{yaniOr}{HTML}{D8A848}
\definecolor{yaniOrClair}{HTML}{E4C15E}
\definecolor{yaniOrProfond}{HTML}{886028}
\definecolor{yaniNoir}{HTML}{1A1712}
\definecolor{yaniCreme}{HTML}{F9F6F0}
\definecolor{yaniGris}{HTML}{6B6252}
\definecolor{yaniGrisClair}{HTML}{EBE3D4}

% ═══════════════════════════════════════════════════════════════════
%  Titres
% ═══════════════════════════════════════════════════════════════════

% Titre de chapitre : un filet doré, le mot « Chapitre N », puis le titre.
\titleformat{\chapter}[display]
  {\normalfont\huge\bfseries\color{yaniNoir}}
  {\filright\normalsize\color{yaniOrProfond}\MakeUppercase{\chaptertitlename\ \thechapter}}
  {8pt}
  {\filright\Large\color{yaniNoir}}
  [\vspace{4pt}{\color{yaniOr}\titlerule[1.5pt]}]

\titlespacing*{\chapter}{0pt}{0pt}{28pt}

% Chapitre NON numéroté : dédicace, remerciements, résumé, listes, table des
% matières, introduction générale, conclusion, webographie. Titre centré,
% souligné du même filet doré que les chapitres numérotés.
%
% C'est titlesec qui doit porter ce style, et non un titre composé à la main
% avant l'appel : \listoffigures, \listoftables et \tableofcontents ouvrent
% eux-mêmes une page par \chapter*, si bien qu'un titre écrit juste avant se
% retrouverait seul sur la page précédente.
\titleformat{name=\chapter,numberless}[display]
  {\normalfont\bfseries\color{yaniNoir}}
  {}
  {0pt}
  {\centering\LARGE}
  [\vspace{6pt}\centering{\color{yaniOr}\rule{0.35\textwidth}{1.5pt}}]

\titlespacing*{name=\chapter,numberless}{0pt}{0pt}{30pt}

\titleformat{\section}
  {\normalfont\large\bfseries\color{yaniOrProfond}}{\thesection}{0.8em}{}
\titleformat{\subsection}
  {\normalfont\normalsize\bfseries\color{yaniNoir}}{\thesubsection}{0.7em}{}
\titleformat{\subsubsection}
  {\normalfont\normalsize\itshape\color{yaniNoir}}{\thesubsubsection}{0.7em}{}

\setcounter{secnumdepth}{3}
\setcounter{tocdepth}{2}

% ═══════════════════════════════════════════════════════════════════
%  En-têtes et pieds de page
% ═══════════════════════════════════════════════════════════════════
\pagestyle{fancy}
\fancyhf{}
\renewcommand{\headrulewidth}{0.4pt}
\renewcommand{\footrulewidth}{0pt}
\fancyhead[L]{\footnotesize\itshape\color{yaniGris}\nouvelleTitreCourt}
\fancyhead[R]{\footnotesize\itshape\color{yaniGris}\leftmark}
\fancyfoot[C]{\thepage}

% `plain` (première page de chapitre) garde le numéro, sans en-tête.
\fancypagestyle{plain}{%
  \fancyhf{}%
  \renewcommand{\headrulewidth}{0pt}%
  \fancyfoot[C]{\thepage}%
}

% \leftmark sans le « CHAPITRE 1. » automatique
\renewcommand{\chaptermark}[1]{\markboth{#1}{}}

% ═══════════════════════════════════════════════════════════════════
%  Numérotation des figures et tableaux par chapitre (Figure 1.1, …)
% ═══════════════════════════════════════════════════════════════════
\counterwithin{figure}{chapter}
\counterwithin{table}{chapter}

% ═══════════════════════════════════════════════════════════════════
%  Figures : image réelle si le fichier existe, cadre d'attente sinon
% ═══════════════════════════════════════════════════════════════════
%
%  C'est le mécanisme qui permet de rédiger et de compiler le rapport AVANT
%  d'avoir produit les captures d'écran et les diagrammes UML. Chaque appel
%  cherche le fichier dans `figures/` :
%
%   • le fichier est là  → l'image est insérée à la largeur demandée ;
%   • le fichier manque  → un cadre gris s'affiche à sa place, portant le nom
%                          de fichier attendu et la légende.
%
%  Il n'y a donc RIEN à modifier dans le texte le jour où les images arrivent :
%  il suffit de déposer `figures/<nom>.png` et de recompiler.
%
%  Usage :  \figProjet[largeur]{nom-du-fichier}{Légende}{cle-de-label}
%           largeur  : fraction de \textwidth (0.85 par défaut)
%           renvoi   : \figref{cle-de-label}
%
\newcommand{\cadreAttente}[2]{%
  \setlength{\fboxrule}{1pt}%
  \setlength{\fboxsep}{0pt}%
  \fcolorbox{yaniOr}{yaniCreme}{%
    \begin{minipage}[c][4.2cm][c]{0.85\textwidth}
      \centering
      {\color{yaniOrProfond}\Large$\square$}\\[6pt]
      {\color{yaniGris}\small\textsc{figure à insérer}}\\[8pt]
      {\color{yaniNoir}\small\textbf{#2}}\\[8pt]
      {\color{yaniGris}\footnotesize\ttfamily figures/#1}
    \end{minipage}%
  }%
}

\newcommand{\figProjet}[4][0.85]{%
  \begin{figure}[H]
    \centering
    \IfFileExists{figures/#2}%
      {\includegraphics[width=#1\textwidth,keepaspectratio]{figures/#2}}%
      {\cadreAttente{#2}{#3}}%
    \caption{#3}
    \label{fig:#4}
  \end{figure}%
}

% Deux figures côte à côte (captures d'écran mobile, souvent étroites).
\newcommand{\figDouble}[6][0.42]{%
  \begin{figure}[H]
    \centering
    \begin{minipage}[t]{0.48\textwidth}
      \centering
      \IfFileExists{figures/#2}%
        {\includegraphics[width=#1\textwidth,keepaspectratio]{figures/#2}}%
        {\cadreAttenteEtroit{#2}{#3}}%
    \end{minipage}\hfill
    \begin{minipage}[t]{0.48\textwidth}
      \centering
      \IfFileExists{figures/#4}%
        {\includegraphics[width=#1\textwidth,keepaspectratio]{figures/#4}}%
        {\cadreAttenteEtroit{#4}{#5}}%
    \end{minipage}
    \caption{#3 · #5}
    \label{fig:#6}
  \end{figure}%
}

\newcommand{\cadreAttenteEtroit}[2]{%
  \setlength{\fboxrule}{1pt}%
  \setlength{\fboxsep}{0pt}%
  \fcolorbox{yaniOr}{yaniCreme}{%
    \begin{minipage}[c][4.2cm][c]{0.92\textwidth}
      \centering
      {\color{yaniOrProfond}\large$\square$}\\[5pt]
      {\color{yaniGris}\scriptsize\textsc{figure à insérer}}\\[6pt]
      {\color{yaniNoir}\scriptsize\textbf{#2}}\\[6pt]
      {\color{yaniGris}\scriptsize\ttfamily figures/#1}
    \end{minipage}%
  }%
}

\newcommand{\figref}[1]{figure~\ref{fig:#1}}
\newcommand{\tabref}[1]{tableau~\ref{tab:#1}}

% ═══════════════════════════════════════════════════════════════════
%  Tableaux — style commun
% ═══════════════════════════════════════════════════════════════════
\newcolumntype{L}[1]{>{\raggedright\arraybackslash}p{#1}}
\newcolumntype{C}[1]{>{\centering\arraybackslash}p{#1}}
\renewcommand{\arraystretch}{1.25}

% En-tête de tableau : fond doré pâle, texte sombre.
\newcommand{\entete}[1]{\cellcolor{yaniGrisClair}\textbf{#1}}

% ═══════════════════════════════════════════════════════════════════
%  Code source
% ═══════════════════════════════════════════════════════════════════
\lstdefinestyle{yani}{
  basicstyle=\ttfamily\scriptsize,
  keywordstyle=\color{yaniOrProfond}\bfseries,
  commentstyle=\color{yaniGris}\itshape,
  stringstyle=\color{yaniOrProfond},
  numbers=none,
  frame=single,
  rulecolor=\color{yaniGrisClair},
  backgroundcolor=\color{yaniCreme},
  breaklines=true,
  showstringspaces=false,
  columns=fullflexible,
  xleftmargin=6pt,
  framexleftmargin=6pt,
  literate=%
    {é}{{\'e}}1 {è}{{\`e}}1 {ê}{{\^e}}1 {à}{{\`a}}1
    {ç}{{\c c}}1 {ù}{{\`u}}1 {û}{{\^u}}1 {î}{{\^i}}1 {ô}{{\^o}}1
    {É}{{\'E}}1 {—}{{---}}1 {’}{{'}}1 {«}{{\guillemotleft}}1 {»}{{\guillemotright}}1
}
\lstset{style=yani}

% ═══════════════════════════════════════════════════════════════════
%  Pages liminaires
% ═══════════════════════════════════════════════════════════════════
%
%  Ouvre une page non numérotée : titre centré (style titlesec ci-dessus),
%  entrée dans la table des matières, et en-tête courant correspondant.
%  L'ancre hyperref est posée AVANT le titre, pour que le lien de la table des
%  matières arrive en haut de la page et non sous le titre.
\newcommand{\pageLiminaire}[1]{%
  \cleardoublepage
  \phantomsection
  \chapter*{#1}%
  \addcontentsline{toc}{chapter}{#1}%
  \markboth{#1}{}%
}

% ═══════════════════════════════════════════════════════════════════
%  Encadré « point d'attention » — utilisé pour les pièges d'exploitation
% ═══════════════════════════════════════════════════════════════════
\newenvironment{attention}[1][Point d'attention]{%
  \par\vspace{6pt}
  \begin{list}{}{%
    \setlength{\leftmargin}{12pt}\setlength{\rightmargin}{12pt}
    \setlength{\topsep}{0pt}}
  \item[]
  \setlength{\fboxrule}{0pt}%
  {\color{yaniOrProfond}\textbf{#1.}}\ %
}{%
  \end{list}\vspace{6pt}\par
}

% ═══════════════════════════════════════════════════════════════════
%  Raccourcis de nommage
% ═══════════════════════════════════════════════════════════════════
\newcommand{\yani}{\textit{Yani Concept by Fati}}

% Valeur restant à renseigner : affichée en gris italique, elle se repère d'un
% coup d'œil dans le PDF et se coupe normalement en fin de ligne — au
% contraire d'une suite de points de suspension, qui déborde des cellules de
% tableau étroites.
\newcommand{\aCompleter}[1]{\textit{\textcolor{yaniGris}{#1}}}


\begin{document}

% ═══════════════════════════════════════════════════════════════════
%  PAGE DE GARDE
% ═══════════════════════════════════════════════════════════════════
%  Toutes les positions sont mesurées depuis le COIN SUPÉRIEUR GAUCHE de la
%  page (`current page.north west`), en centimètres, y négatif vers le bas.
%  Pour descendre un bloc, augmenter son deuxième nombre.
%
%  Repères verticaux : logos −2,3 · filière −6,6 · cadre du titre −7,9 à −12,9
%  · « Réalisé par » −16,5 · « Encadré par » −20,6 · année −27,6.
\begin{titlepage}
\thispagestyle{empty}
\begin{tikzpicture}[remember picture, overlay]

  % ---- Bande jaune en haut à gauche ----
  \fill[jauneEMSI] ($(current page.north west)+(0.4cm,-1.1cm)$)
    rectangle ($(current page.north west)+(8cm,-1.5cm)$);

  % ---- Logos : l'école à gauche, le projet à droite ----
  \node[anchor=north west] at ($(current page.north west)+(0.7cm,-2.3cm)$) {%
    \logoCouverture{EMSI.png}{2.1cm}
  };
  \node[anchor=north east] at ($(current page.north east)+(-1cm,-2.3cm)$) {%
    \logoCouverture{logo.png}{2.3cm}
  };

  % ---- Bande jaune sous les logos ----
  \fill[jauneEMSI] ($(current page.north west)+(3.6cm,-5.3cm)$)
    rectangle ($(current page.north east)+(-0.4cm,-5.7cm)$);

  % ---- Nature du travail + filière ----
  \node[anchor=north east] at ($(current page.north east)+(-1cm,-6.6cm)$) {%
    \begin{minipage}{9cm}\raggedleft
      {\large\bfseries \nouvelleNature}\\[4pt]
      {\normalsize\bfseries\color{vertEMSI} \nouvelleFiliere}
    \end{minipage}
  };

  % ---- Ombre + cadre arrondi du titre, avec ses deux attaches jaunes ----
  \fill[ombre, rounded corners=8pt] ($(current page.north west)+(0.9cm,-8.1cm)$)
    rectangle ($(current page.north west)+(11.2cm,-13.1cm)$);
  \draw[jauneEMSI, line width=3pt, rounded corners=8pt]
    ($(current page.north west)+(0.7cm,-7.9cm)$)
    rectangle ($(current page.north west)+(11cm,-12.9cm)$);
  \fill[jauneEMSI] ($(current page.north west)+(11cm,-8.3cm)$)
    rectangle ($(current page.north east)+(-0.4cm,-8.7cm)$);
  \fill[jauneEMSI] ($(current page.north west)+(11cm,-12.1cm)$)
    rectangle ($(current page.north west)+(11.4cm,-12.9cm)$);
  % La césure est désactivée ici seulement : un titre coupé en « ap-plication »
  % au milieu d'une couverture se voit tout de suite. Le reste du rapport, lui,
  % garde la césure — sans elle, le texte justifié se troue.
  %
  % \hyphenchar et non \hyphenpenalty : dans un nœud TikZ, la pénalité est
  % rétablie avant que le paragraphe ne soit coupé, donc sans effet. Retirer le
  % caractère de césure de la police, après l'avoir sélectionnée, tient.
  \node[anchor=center, align=center, text width=9.6cm]
    at ($(current page.north west)+(5.85cm,-10.4cm)$) {%
    \normalsize\bfseries\selectfont
    \hyphenchar\font=-1\relax
    \nouvelleTitre
  };

  % ---- Filet jaune de séparation ----
  \fill[jauneEMSI] ($(current page.north west)+(1cm,-13.9cm)$)
    rectangle ($(current page.north east)+(-1cm,-14.05cm)$);

  % ---- Bloc « Réalisé par » ----
  % La barre est calée sur UN auteur. Pour un binôme, ajouter la seconde ligne
  % dans le minipage ci-dessous et descendre le −17.9 à −18.5.
  \fill[jauneEMSI] ($(current page.north west)+(1.1cm,-16.5cm)$)
    rectangle ($(current page.north west)+(1.4cm,-17.9cm)$);
  \node[anchor=north west] at ($(current page.north west)+(1.7cm,-16.6cm)$) {%
    \begin{minipage}{12cm}
      {\bfseries\color{vertEMSI} Réalisé par~:}\\[6pt]
      \nouvelleAuteur
    \end{minipage}
  };

  % ---- Bloc « Encadré par » ----
  \fill[jauneEMSI] ($(current page.north west)+(1.1cm,-20.6cm)$)
    rectangle ($(current page.north west)+(1.4cm,-23.4cm)$);
  \node[anchor=north west] at ($(current page.north west)+(1.7cm,-20.7cm)$) {%
    \begin{minipage}{15.5cm}
      {\bfseries\color{vertEMSI} Encadré par~:}\\[6pt]
      \nouvelleEncadrantPedagogique, Encadrant Académique,
      EMSI \nouvelleVille.\\[6pt]
      \nouvelleEncadrantEntreprise, Encadrant Professionnel,
      \nouvelleTitreCourt, \nouvelleVille.
    \end{minipage}
  };

  % ---- Filet jaune de bas de page ----
  \fill[jauneEMSI] ($(current page.north west)+(1cm,-26.7cm)$)
    rectangle ($(current page.north east)+(-1cm,-26.85cm)$);

  % ---- Année universitaire ----
  \node[anchor=north] at ($(current page.north)+(0cm,-27.6cm)$) {%
    {\large\bfseries Année Universitaire~: \nouvelleAnnee}
  };

\end{tikzpicture}
\end{titlepage}

% ═══════════════════════════════════════════════════════════════════
%  PAGES LIMINAIRES — numérotation romaine
% ═══════════════════════════════════════════════════════════════════
\pagenumbering{roman}
\setcounter{page}{1}

% ═══════════════════════════════════════════════════════════════════
%  Pages liminaires : dédicace, remerciements, résumé, abstract, glossaire
% ═══════════════════════════════════════════════════════════════════

% ── Dédicace ─────────────────────────────────────────────────────
\pageLiminaire{Dédicace}

\vspace{1.5cm}
\begin{center}
\begin{minipage}{0.82\textwidth}
\itshape
\large

Je dédie ce modeste travail :

\vspace{0.7cm}

À mes parents, pour leurs sacrifices, leur patience et la confiance qu'ils
n'ont jamais cessé de m'accorder.

\vspace{0.5cm}

À ma famille et à mes amis, pour leur soutien constant et leurs
encouragements tout au long de ces années.

\vspace{0.5cm}

À tous les professeurs de \nouvelleEtablissement, qui ont donné sans compter
de leur temps et de leur savoir pour nous former.

\vspace{0.5cm}

À Madame Fati, gérante du centre \yani, pour m'avoir ouvert les portes de son
institut, pour le temps qu'elle a consacré à répondre à mes questions et pour
la confiance qu'elle a placée dans ce projet.

\vspace{0.5cm}

À toutes celles et ceux qui, de près ou de loin, ont participé à
l'élaboration de ce travail.

\vspace{0.5cm}

Et à toute personne qui prendra ce rapport entre ses mains.

\end{minipage}
\end{center}

% ── Remerciements ────────────────────────────────────────────────
\pageLiminaire{Remerciements}

C'est avec une grande satisfaction que je tiens à remercier toutes les
personnes qui m'ont aidé, accompagné et soutenu jusqu'à l'aboutissement de ce
projet de fin d'études.

\vspace{0.3cm}

J'adresse en premier lieu mes remerciements à l'ensemble du corps professoral
de \nouvelleEtablissement, dont les enseignements constituent le socle sur
lequel ce travail a pu être bâti. Les notions de génie logiciel, de bases de
données, de sécurité et de conduite de projet mobilisées dans les pages qui
suivent y ont toutes été acquises.

\vspace{0.3cm}

Je remercie tout particulièrement \nouvelleEncadrantPedagogique, mon
encadrant pédagogique, pour son suivi régulier, ses relectures exigeantes et
les conseils qui ont orienté ce projet aux moments où il en avait le plus
besoin.

\vspace{0.3cm}

Je remercie également \nouvelleEncadrantEntreprise, ainsi que Madame Fati,
gérante du centre \yani, pour sa disponibilité et pour la qualité des échanges
que nous avons eus. C'est en observant le fonctionnement réel de son institut
--- le carnet de rendez-vous, les appels téléphoniques, la carte de fidélité
tamponnée --- que les besoins de cette application ont pu être formulés
autrement que de façon théorique. Les allers-retours réguliers avec elle ont
directement façonné les arbitrages de conception exposés dans ce rapport :
c'est à sa demande, par exemple, que la durée d'une prestation a cessé de
commander la grille des créneaux.

\vspace{0.3cm}

Mes remerciements vont enfin aux membres du jury, qui me font l'honneur
d'accepter d'évaluer ce travail, ainsi qu'à mes parents, dont le soutien
indéfectible ne s'est jamais démenti.

\vspace{0.5cm}

Merci à toutes et à tous.

% ── Résumé ───────────────────────────────────────────────────────
\pageLiminaire{Résumé}

Ce rapport de projet de fin d'études présente la conception et la réalisation
d'une solution logicielle complète pour \yani, un centre de beauté situé à
Casablanca. Le projet répond à une situation courante dans les petites
structures de services : une activité entièrement pilotée à la main --- carnet
de rendez-vous papier, prises de rendez-vous par téléphone et par messagerie
instantanée, vente de produits sans suivi de stock, fidélité gérée par une
carte tamponnée --- avec les conséquences qui en découlent : double
réservation d'un même créneau, absence d'historique exploitable, temps
consacré au téléphone au détriment des clientes présentes.

\vspace{0.3cm}

La solution développée est un mono-dépôt réunissant trois applications
adossées à une même base de données PostgreSQL. Une \textbf{API REST} bâtie
avec NestJS et Prisma porte l'ensemble des règles métier. Une
\textbf{application mobile} développée avec React Native et Expo permet aux
clientes de consulter les catalogues, de réserver une prestation, de commander
des produits et de suivre leur compte de fidélité. Un \textbf{back-office} web
en React donne à la gérante et à son personnel la maîtrise des commandes, des
rendez-vous, du catalogue, des horaires d'ouverture et du programme de
fidélité.

\vspace{0.3cm}

Le rapport détaille l'ensemble du cycle de développement : étude de
l'existant, expression des besoins, modélisation UML, choix techniques,
architecture en couches, implémentation et validation. Une méthode agile
inspirée de SCRUM, organisée en six sprints, a permis de livrer par
incréments et de confronter chaque version au fonctionnement réel de
l'institut. Une attention particulière a été portée aux propriétés que les
tests unitaires classiques ne savent pas démontrer : la sérialisation des
réservations concurrentes par verrou consultatif PostgreSQL, l'atomicité du
crédit de fidélité, l'idempotence de l'émission des bons de récompense et
l'anonymisation des comptes supprimés, exigée par la loi 09-08 sur la
protection des données personnelles. Vingt-cinq suites de tests, dont huit
écrivent dans une véritable base PostgreSQL, sont rejouées à chaque
modification par une chaîne d'intégration continue.

\vspace{0.3cm}

Le résultat est une plate-forme fonctionnelle, trilingue (français, arabe,
anglais), documentée et prête au déploiement, qui supprime le carnet papier,
rend la réservation possible à toute heure et donne à l'institut une visibilité
qu'il n'avait pas sur son activité.

\vspace{0.6cm}

\textbf{Mots clés :} application mobile, API REST, NestJS, Prisma, PostgreSQL,
React Native, Expo, React, TypeScript, réservation en ligne, programme de
fidélité, back-office, concurrence, sécurité applicative, internationalisation,
intégration continue.

% ── Abstract ─────────────────────────────────────────────────────
\pageLiminaire{Abstract}

\selectlanguage{english}

This final year project report presents the design and implementation of a
complete software solution for \textit{Yani Concept by Fati}, a beauty centre
located in Casablanca. The project addresses a situation common to small
service businesses: an activity managed entirely by hand --- a paper
appointment book, bookings taken over the phone and through instant messaging,
product sales with no stock tracking, and a loyalty programme run on a stamped
card --- with the consequences that follow: the same slot booked twice, no
usable history, and time spent on the phone at the expense of the customers
already in the salon.

\vspace{0.3cm}

The delivered solution is a monorepo bringing together three applications
backed by a single PostgreSQL database. A \textbf{REST API} built with NestJS
and Prisma carries all the business rules. A \textbf{mobile application} built
with React Native and Expo lets customers browse the catalogues, book a
treatment, order products and follow their loyalty account. A \textbf{web
back-office} written in React gives the manager and her staff control over
orders, appointments, the catalogue, opening hours and the loyalty programme.

\vspace{0.3cm}

The report covers the full development cycle: study of the existing process,
requirements elicitation, UML modelling, technical choices, layered
architecture, implementation and validation. An agile method inspired by SCRUM,
organised into six sprints, made it possible to deliver in increments and to
confront each version with the real workings of the salon. Particular care was
given to properties that ordinary unit tests cannot demonstrate: serialisation
of concurrent bookings through a PostgreSQL advisory lock, atomicity of loyalty
crediting, idempotent issuance of reward vouchers, and anonymisation of deleted
accounts as required by Moroccan law 09-08 on personal data protection.
Twenty-five test suites, eight of which write to a real PostgreSQL database,
are replayed on every change by a continuous integration pipeline.

\vspace{0.3cm}

The result is a working, trilingual (French, Arabic, English), documented
platform, ready for deployment, which removes the paper book, makes booking
possible at any hour, and gives the salon a visibility over its own activity
that it previously lacked.

\vspace{0.6cm}

\textbf{Keywords:} mobile application, REST API, NestJS, Prisma, PostgreSQL,
React Native, Expo, React, TypeScript, online booking, loyalty programme,
back-office, concurrency, application security, internationalisation,
continuous integration.

\selectlanguage{french}

% ── Glossaire ────────────────────────────────────────────────────
\pageLiminaire{Glossaire}

\begin{center}
\begin{tabularx}{\textwidth}{L{2.6cm} X}
\toprule
\entete{Sigle} & \entete{Signification} \\
\midrule
API      & \textit{Application Programming Interface} --- interface de programmation applicative \\
Argon2   & Fonction de dérivation de clé, lauréate du \textit{Password Hashing Competition} \\
CI       & \textit{Continuous Integration} --- intégration continue \\
CNDP     & Commission Nationale de contrôle de la protection des Données à caractère Personnel \\
COD      & \textit{Cash On Delivery} --- paiement à la remise ou à la livraison \\
CORS     & \textit{Cross-Origin Resource Sharing} --- partage de ressources entre origines \\
CSP      & \textit{Content Security Policy} --- politique de sécurité du contenu \\
DTO      & \textit{Data Transfer Object} --- objet de transfert de données \\
EAS      & \textit{Expo Application Services} --- chaîne de compilation native d'Expo \\
HSTS     & \textit{HTTP Strict Transport Security} \\
HTTP(S)  & \textit{HyperText Transfer Protocol (Secure)} \\
IHM      & Interface Homme--Machine \\
i18n     & \textit{Internationalization} --- internationalisation \\
IANA     & \textit{Internet Assigned Numbers Authority} --- source de la base des fuseaux horaires \\
JSON     & \textit{JavaScript Object Notation} \\
JWT      & \textit{JSON Web Token} --- jeton d'authentification signé \\
ORM      & \textit{Object-Relational Mapping} --- correspondance objet-relationnel \\
PFE      & Projet de Fin d'Études \\
REST     & \textit{REpresentational State Transfer} --- style d'architecture d'API \\
RGPD     & Règlement Général sur la Protection des Données \\
SGBD     & Système de Gestion de Bases de Données \\
SMTP     & \textit{Simple Mail Transfer Protocol} \\
SPA      & \textit{Single Page Application} --- application web mono-page \\
SQL      & \textit{Structured Query Language} \\
TDD      & \textit{Test-Driven Development} --- développement piloté par les tests \\
TTL      & \textit{Time To Live} --- durée de validité \\
UI / UX  & \textit{User Interface} / \textit{User eXperience} \\
UML      & \textit{Unified Modeling Language} --- langage de modélisation unifié \\
UTC      & \textit{Coordinated Universal Time} --- temps universel coordonné \\
UUID     & \textit{Universally Unique Identifier} \\
XSS      & \textit{Cross-Site Scripting} \\
\bottomrule
\end{tabularx}
\end{center}


% ── Listes automatiques ──────────────────────────────────────────
% Ces trois commandes ouvrent elles-mêmes leur page et composent leur titre
% (style « chapitre non numéroté » défini dans le préambule) ; il ne reste
% qu'à poser l'ancre et l'entrée de table des matières juste avant.
\cleardoublepage
\phantomsection
\addcontentsline{toc}{chapter}{Liste des figures}
\markboth{Liste des figures}{}
\listoffigures

\cleardoublepage
\phantomsection
\addcontentsline{toc}{chapter}{Liste des tableaux}
\markboth{Liste des tableaux}{}
\listoftables

\cleardoublepage
\phantomsection
\addcontentsline{toc}{chapter}{Table des matières}
\markboth{Table des matières}{}
\tableofcontents

% ═══════════════════════════════════════════════════════════════════
%  CORPS DU RAPPORT — numérotation arabe
% ═══════════════════════════════════════════════════════════════════
\cleardoublepage
\pagenumbering{arabic}
\setcounter{page}{1}

% ═══════════════════════════════════════════════════════════════════
%  Introduction générale
% ═══════════════════════════════════════════════════════════════════
\pageLiminaire{Introduction Générale}

La transformation numérique n'est plus l'apanage des grandes entreprises. Elle
touche aujourd'hui les commerces de proximité, les cabinets, les artisans et
les instituts, dont l'activité repose encore très souvent sur des outils
manuels : un carnet, un téléphone, une carte cartonnée. Ces outils ont une
qualité que l'on sous-estime --- ils sont immédiats, personne n'a besoin de les
apprendre --- et un défaut qui devient handicapant dès que l'activité croît :
ils ne conservent rien d'exploitable, ne se consultent qu'à un seul endroit, et
n'existent que pendant les heures d'ouverture.

Le secteur de la beauté et du bien-être illustre ce constat de façon
particulièrement nette. L'activité d'un institut s'organise autour d'une
ressource rare et non stockable --- un créneau de cabine --- que plusieurs
canaux se disputent en même temps : la cliente qui appelle, celle qui écrit sur
une messagerie instantanée, et celle qui se présente au comptoir. Le carnet
papier arbitre entre ces canaux, mais il ne peut le faire que si quelqu'un le
tient ouvert devant soi. Toute réservation prise ailleurs qu'à côté du carnet
est une réservation à risque.

C'est dans ce contexte que s'inscrit ce projet de fin d'études. \yani\ est un
centre de beauté installé à Casablanca, qui propose des prestations de soin et
vend des produits cosmétiques. Sa gérante y assure à la fois les soins, la
prise des rendez-vous, le suivi du stock et la fidélisation de sa clientèle.
L'objectif du projet est de porter l'ensemble de ces activités sur une
plate-forme numérique : une application mobile pour les clientes, un
back-office web pour l'institut, et une API qui porte les règles métier
communes aux deux.

Le travail présenté ici ne s'est pas limité à traduire un carnet papier en
formulaires. Une réservation en ligne pose des questions qu'un carnet ne pose
pas : que se passe-t-il lorsque deux clientes valident le même créneau à la
même seconde ? Comment un institut ouvert de 9\,h à 13\,h puis de 14\,h à
18\,h, à l'heure de Casablanca, expose-t-il ses disponibilités à une
application qui raisonne en temps universel ? Que devient l'historique
d'achat d'une cliente qui exerce son droit à l'effacement, alors même que cet
historique porte le chiffre d'affaires de l'institut ? Ces questions ont
structuré la conception autant que les fonctionnalités elles-mêmes, et elles
occupent une place proportionnée dans ce rapport.

Le présent document rend compte de ce travail. Il s'articule autour de quatre
chapitres.

Le \textbf{premier chapitre} présente l'organisme d'accueil, l'étude de
l'existant, la problématique qui en découle et la solution retenue, puis la
méthode de conduite de projet et sa planification.

Le \textbf{deuxième chapitre} est consacré à l'analyse et à la conception : la
spécification des besoins fonctionnels et non fonctionnels, les acteurs du
système, puis la modélisation UML --- diagrammes de cas d'utilisation,
diagrammes de séquence et diagramme de classes.

Le \textbf{troisième chapitre} expose l'étude technique : les langages, les
cadriciels et le système de gestion de base de données retenus, les
environnements de travail, et l'architecture technique de la solution, y
compris la chaîne de sécurité appliquée à chaque requête.

Le \textbf{quatrième chapitre} traite de la mise en \oe uvre : la structure
effective du code des trois applications, la stratégie de tests et
l'intégration continue, la présentation des interfaces réalisées, et enfin le
déploiement et l'exploitation de la plate-forme.

Une conclusion générale dresse le bilan du travail accompli, ses limites
assumées et les perspectives d'évolution de la solution.

% ═══════════════════════════════════════════════════════════════════
%  Chapitre 1 — Présentation de l'organisme d'accueil et du projet
% ═══════════════════════════════════════════════════════════════════
\chapter{Présentation de l'organisme d'accueil et du contexte du projet}

\section*{Introduction}
\addcontentsline{toc}{section}{Introduction}

Ce chapitre présente le cadre dans lequel ce projet de fin d'études a été
réalisé. Il décrit d'abord l'organisme d'accueil --- son activité, son
organisation et les personnes qui y travaillent --- puis le contexte du projet
proprement dit : la façon dont l'institut fonctionnait avant l'intervention,
les difficultés que ce fonctionnement engendrait, la solution retenue pour y
répondre, ainsi que les enjeux et les contraintes qui l'ont encadrée. Il se
termine par la méthode de conduite de projet adoptée et par la planification
qui en découle.

% ─────────────────────────────────────────────────────────────────
\section{Présentation de l'organisme d'accueil}

\subsection{Le centre \yani}

\yani\ est un centre de beauté et de bien-être installé à Marrakech. C'est une
structure de proximité : une clientèle essentiellement féminine, largement
fidélisée, qui revient à intervalles réguliers et qui connaît personnellement
la gérante. Cette proximité est le principal atout commercial de l'institut ;
c'est aussi ce qui rend la question de l'outil informatique délicate, car toute
solution qui interposerait une machine froide entre la cliente et l'institut
détruirait précisément ce qui fait sa valeur.

\figProjet[0.32]{logo.png}{Identité visuelle de \yani}{logo}

L'identité visuelle du centre --- un or lumineux sur fond crème ou noir profond
--- a été reprise à l'identique dans les trois applications développées. Ce
n'est pas un détail cosmétique : pour une cliente, l'application doit
manifestement être celle de l'institut qu'elle connaît, et non un outil
générique de réservation dans lequel l'institut ne serait qu'une entrée parmi
d'autres.

\subsection{Activités et offre de services}

L'activité du centre se répartit en deux volets, qui ont chacun leur logique
propre et qui sont tous deux couverts par la solution développée.

\begin{itemize}[leftmargin=1.2cm]
  \item \textbf{Les prestations de soin.} Elles constituent le c\oe ur de
  métier. Chaque prestation appartient à une catégorie (soins du visage, soins
  du corps, épilation, coiffure, onglerie\ldots), porte un tarif et
  s'effectue en cabine, sur rendez-vous. La ressource limitante est le nombre
  de cabines réservables simultanément.

  \item \textbf{La vente de produits cosmétiques.} L'institut revend des
  produits de soin. Cette activité obéit à une logique différente : elle porte
  sur un stock physique, se règle à la remise ou à la livraison, et n'occupe
  aucune cabine.
\end{itemize}

À ces deux volets s'ajoute un \textbf{programme de fidélité}, transversal, qui
récompense à la fois les prestations réalisées et les produits achetés.

\subsection{Organisation et acteurs}

L'institut est une structure de petite taille, ce qui a une conséquence directe
sur la conception : les rôles ne sont pas cloisonnés comme ils le seraient dans
une entreprise plus grande, et une même personne assure souvent plusieurs
fonctions dans la même journée. Trois profils ont néanmoins été distingués,
parce qu'ils correspondent à trois niveaux de responsabilité réellement
différents.

\begin{table}[H]
\centering
\caption{Les acteurs de l'institut et leurs responsabilités}
\label{tab:acteurs-institut}
\begin{tabularx}{\textwidth}{L{3.1cm} L{2.4cm} X}
\toprule
\entete{Acteur} & \entete{Rôle système} & \entete{Responsabilités} \\
\midrule
La gérante & \texttt{ADMIN} &
Direction du centre. Fixe les tarifs et le catalogue, règle les horaires
d'ouverture et la capacité, définit le programme de fidélité, gère les comptes
du personnel et consulte les exports comptables. \\

Le personnel & \texttt{STAFF} &
Réalise les soins et tient le comptoir. Traite les commandes, confirme et
clôture les rendez-vous, honore les bons de récompense, crédite
exceptionnellement des points à la main. \\

Les clientes & \texttt{CLIENT} &
Consultent les catalogues, réservent leurs rendez-vous, commandent des
produits et suivent leur compte de fidélité. \\
\bottomrule
\end{tabularx}
\end{table}

\figProjet[0.75]{organisation-centre.png}{Organisation du centre et acteurs du système}{organisation}

Cette hiérarchie de rôles n'est pas une simple convention documentaire : elle
est appliquée par le serveur sur chaque route de l'API, indépendamment de ce
que l'interface propose ou masque. Une cliente qui obtiendrait par un moyen
quelconque l'adresse d'une route d'administration se verrait refuser l'accès
par le serveur lui-même. Ce point est développé au chapitre~4.

% ─────────────────────────────────────────────────────────────────
\section{Présentation du contexte du projet}

\subsection{Vue d'ensemble du projet}

Le projet consiste à concevoir et à réaliser une solution logicielle complète
couvrant l'ensemble de l'activité de l'institut. Elle prend la forme d'un
mono-dépôt réunissant trois applications qui partagent une seule base de
données :

\begin{itemize}[leftmargin=1.2cm]
  \item une \textbf{API REST} qui porte les données et la totalité des règles
  métier ;
  \item une \textbf{application mobile} destinée aux clientes ;
  \item un \textbf{back-office web} destiné à la gérante et à son personnel.
\end{itemize}

Le choix de faire porter toutes les règles métier par l'API, et par elle seule,
est structurant. Les deux interfaces --- mobile et web --- n'implémentent
aucune règle : elles affichent ce que le serveur leur donne et lui soumettent
ce que l'utilisatrice saisit. Une règle telle que « on ne peut pas réserver
dans le passé » ou « un rendez-vous annulé ne peut plus être marqué terminé »
n'existe qu'à un seul endroit. Ce principe évite la divergence classique où
l'application mobile et le back-office finissent par appliquer des règles
subtilement différentes à la même donnée.

\subsection{Étude de l'existant}

Avant tout développement, une observation du fonctionnement réel de l'institut
a été menée. Elle a mis en évidence quatre processus manuels, tous fonctionnels
au quotidien, et tous porteurs des mêmes limites.

\paragraph{La prise de rendez-vous.} Elle s'effectue par trois canaux
concurrents : le téléphone, la messagerie instantanée et la présentation
directe au comptoir. La réservation est inscrite au crayon dans un carnet
papier, unique et physiquement situé à l'accueil. La cliente ne dispose
d'aucune confirmation écrite autre que la mémoire de l'échange.

\paragraph{La vente de produits.} Elle se fait exclusivement sur place. Le
stock n'est suivi par aucun outil : la gérante l'estime de visu. Un produit
épuisé ne se découvre qu'au moment où une cliente le demande.

\paragraph{La fidélité.} Elle repose sur une carte cartonnée tamponnée à chaque
visite. La carte perdue emporte l'historique avec elle ; la carte oubliée à la
maison prive la cliente de son tampon ; et l'institut n'a aucune vision
consolidée de ce que le programme lui coûte ni de ce qu'il lui rapporte.

\paragraph{Le suivi de l'activité.} Il n'existe pas de manière systématique. Il
n'est pas possible de répondre, sans reconstitution manuelle, à des questions
aussi élémentaires que « combien de rendez-vous ce mois-ci ? », « quelles sont
les prestations les plus demandées ? » ou « quelles clientes ne sont pas
revenues depuis six mois ? ».

\subsection{Problématique}

Ce fonctionnement, viable tant que l'activité reste modeste, engendre des
difficultés qui se renforcent mutuellement à mesure qu'elle croît.

\begin{itemize}[leftmargin=1.2cm]
  \item \textbf{Le risque de double réservation.} Trois canaux alimentent un
  carnet unique. Une réservation prise par téléphone pendant qu'une cliente
  écrit sur la messagerie peut viser le même créneau. Le conflit ne se découvre
  qu'au moment où les deux clientes se présentent, et sa résolution se fait
  toujours au détriment de l'une d'elles.

  \item \textbf{L'indisponibilité hors des heures d'ouverture.} Une cliente qui
  souhaite réserver le soir ou le dimanche doit attendre. Une partie de ces
  demandes ne se transforme jamais en rendez-vous.

  \item \textbf{Le temps consacré au téléphone.} Chaque appel interrompt un
  soin en cours. Ce coût est invisible dans les comptes, mais il est réel : il
  se paie en qualité de service pour la cliente présente en cabine.

  \item \textbf{L'absence de mémoire exploitable.} Le carnet et la carte de
  fidélité ne conservent qu'une trace éphémère et non consolidée. L'institut ne
  peut donc ni mesurer son activité, ni identifier ses prestations phares, ni
  détecter l'érosion de sa clientèle.

  \item \textbf{La fragilité du support physique.} Un carnet se perd, s'abîme,
  et n'existe qu'en un seul exemplaire, à un seul endroit.

  \item \textbf{La rupture de stock invisible.} Sans suivi, un produit épuisé
  n'est découvert qu'au comptoir, devant la cliente.
\end{itemize}

Ces difficultés ont un dénominateur commun : \emph{l'information de l'institut
n'existe qu'en un seul exemplaire, dans un seul lieu, et pendant les seules
heures d'ouverture}. C'est cette contrainte que le projet lève.

\subsection{Solution proposée}

La solution retenue est une plate-forme composée de trois applications adossées
à une base de données unique. Ce choix mérite d'être justifié : il aurait été
possible de se contenter d'un back-office web, ou de recourir à une plate-forme
de réservation existante. Ni l'une ni l'autre de ces options n'a été retenue.

Un back-office seul aurait déplacé le carnet papier vers un écran, sans rien
changer pour la cliente : elle aurait continué d'appeler. Une plate-forme de
réservation généraliste, quant à elle, aurait réglé la question des créneaux
mais placé l'institut au milieu d'un annuaire de concurrents, sans couvrir ni
la vente de produits ni un programme de fidélité propre --- et en faisant
dépendre l'activité d'un tiers sur lequel l'institut n'a aucune prise.

La plate-forme développée comprend donc :

\begin{itemize}[leftmargin=1.2cm]
  \item \textbf{Une API REST (NestJS, Prisma, PostgreSQL).} Elle expose plus de
  quatre-vingts routes et porte l'ensemble des règles métier : catalogue,
  moteur de créneaux, machine à états des rendez-vous et des commandes,
  programme de fidélité, authentification, exports.

  \item \textbf{Une application mobile (React Native, Expo).} Vingt-quatre
  écrans permettent à la cliente de parcourir les catalogues, de réserver, de
  commander, de suivre ses points et de gérer son compte. Elle est disponible
  en français, en arabe et en anglais.

  \item \textbf{Un back-office web (React, Vite).} Huit pages donnent à
  l'institut la maîtrise complète de son activité : tableau de bord,
  commandes, rendez-vous, catalogue, fidélité, horaires, utilisateurs, avec
  export Excel des tableaux réservé à la gérante.
\end{itemize}

\begin{table}[H]
\centering
\caption{Réponse apportée à chaque difficulté identifiée}
\label{tab:probleme-solution}
\begin{tabularx}{\textwidth}{L{5.2cm} X}
\toprule
\entete{Difficulté} & \entete{Réponse apportée} \\
\midrule
Double réservation d'un créneau &
Canal de réservation unique, et sérialisation des réservations concurrentes par
verrou consultatif PostgreSQL : deux demandes simultanées sur le même créneau
sont traitées l'une après l'autre, jamais en parallèle. \\

Indisponibilité hors ouverture &
Réservation possible à toute heure depuis l'application, sur la grille de
créneaux calculée par le serveur à partir des horaires réels du centre. \\

Temps passé au téléphone &
La réservation en autonomie décharge l'accueil ; le rendez-vous pris par
téléphone reste possible, le personnel le saisissant pour la cliente. \\

Absence de mémoire &
Historique complet en base, tableau de bord, et exports Excel horodatés des
clientes, des commandes et des rendez-vous. \\

Fragilité du support papier &
Base PostgreSQL, script de sauvegarde horodatée avec rotation à quatorze jours
et procédure de restauration documentée. \\

Rupture de stock invisible &
Stock suivi par produit, contrôlé à la commande, décrémenté de façon atomique à
la confirmation, et alerte de stock faible sur le tableau de bord. \\

Carte de fidélité perdue &
Compte de fidélité attaché au compte de la cliente : solde, historique, paliers
de visites et bons de récompense, consultables depuis l'application. \\
\bottomrule
\end{tabularx}
\end{table}

\subsection{Enjeux du projet}

\begin{itemize}[leftmargin=1.2cm]
  \item \textbf{Enjeu commercial.} Rendre la réservation possible en dehors des
  heures d'ouverture, et transformer en rendez-vous des demandes aujourd'hui
  perdues faute d'interlocuteur disponible.

  \item \textbf{Enjeu de fidélisation.} Remplacer une carte cartonnée fragile
  par un compte consultable, doté de paliers de visites et de bons de
  récompense traçables --- un programme que l'institut peut enfin piloter et
  chiffrer.

  \item \textbf{Enjeu de pilotage.} Donner à la gérante une visibilité qu'elle
  n'a jamais eue sur son activité : commandes en attente, rendez-vous du jour,
  stocks faibles, historique exportable.

  \item \textbf{Enjeu d'image.} Une application aux couleurs de l'institut,
  soignée et trilingue, participe directement à la perception de la marque
  auprès d'une clientèle jeune et connectée.

  \item \textbf{Enjeu de conformité.} Traiter des données personnelles impose
  des obligations, notamment le droit à l'effacement prévu par la loi 09-08 et
  contrôlé par la CNDP. Ce droit devait être réellement exerçable, et non
  simplement annoncé.
\end{itemize}

\subsection{Contraintes}

\paragraph{Contrainte d'exactitude sur les créneaux.} Un créneau de cabine ne
peut pas être vendu deux fois. Cette contrainte, banale à énoncer, est la plus
difficile du projet : la vérification de disponibilité et l'insertion du
rendez-vous portent sur un \emph{ensemble} de lignes, ce qu'aucune écriture
conditionnelle sur une ligne unique ne sait rendre atomique.

\paragraph{Contrainte de fuseau horaire.} L'institut vit à l'heure de
Marrakech, dont le décalage varie dans l'année. Les instants sont stockés en
UTC, les horaires d'ouverture exprimés en heure locale. Toute confusion entre
les deux se traduit par une grille de créneaux décalée d'une heure --- une
erreur silencieuse, qui ne lève aucune exception et qui ne se voit qu'en
comparant l'écran au carnet.

À noter, car cela surprend à la lecture du code : le fuseau est déclaré sous le
nom \texttt{Africa/Casablanca}. Ce n'est pas une erreur ni un reliquat. La base
IANA nomme chaque zone d'après sa ville la plus peuplée, et cette zone couvre
l'ensemble du Maroc, Marrakech comprise. La renommer serait la casser.

\paragraph{Contrainte de conservation de l'historique.} Une cliente peut
demander la suppression de son compte. Ses commandes et ses rendez-vous
portent pourtant le chiffre d'affaires de l'institut, qui doit lui survivre.
Concilier les deux exige une anonymisation, et non une suppression en cascade.

\paragraph{Contrainte d'usage.} L'utilisatrice principale du back-office n'est
pas informaticienne. L'outil devait être immédiatement compréhensible, sans
formation ni documentation à consulter --- ce qui a proscrit toute exposition
de valeurs techniques brutes dans l'interface comme dans les exports.

\paragraph{Contrainte de coût.} L'institut est une petite structure.
L'ensemble de la solution repose donc sur des technologies libres et sur une
infrastructure minimale, déployable sur un seul serveur.

\paragraph{Contrainte de délai.} Le projet devait aboutir à une plate-forme
fonctionnelle et déployable dans la durée impartie au stage de fin d'études,
soit environ huit semaines de développement effectif.

% ─────────────────────────────────────────────────────────────────
\section{Conduite du projet}

\subsection{La méthode agile SCRUM}

Le choix d'une méthode de développement conditionne l'organisation de tout le
projet. La méthode agile SCRUM a été retenue, pour une raison qui tient au
contexte plus qu'à la mode : les besoins de l'institut n'étaient pas
formalisables à l'avance. Ils se sont précisés à mesure que des versions
concrètes étaient montrées à la gérante.

L'exemple le plus net concerne le moteur de créneaux. La première version
calculait les disponibilités à partir de la durée de chaque prestation : une
prestation d'une heure vingt-cinq occupait le temps correspondant. Le
comportement était techniquement juste, et fonctionnellement inutilisable ---
il produisait des heures de rendez-vous irrégulières, impossibles à annoncer au
téléphone. La gérante espace en réalité ses rendez-vous d'un écart fixe, quelle
que soit la prestation. Le moteur a donc été refait : un rendez-vous occupe
exactement un créneau, et la durée de la prestation n'est plus qu'une
information de gestion. Aucune spécification écrite à l'avance n'aurait produit
cette règle ; seule la confrontation d'une version fonctionnelle au métier réel
l'a fait apparaître.

\figProjet[0.72]{scrum.png}{Le cycle de la méthode agile SCRUM}{scrum}

\subsection{Répartition des rôles}

Le projet ayant été mené par un seul développeur, la répartition canonique des
rôles SCRUM a été adaptée sans être abandonnée --- car ce sont bien trois
responsabilités distinctes, même lorsque deux d'entre elles sont portées par la
même personne.

\begin{table}[H]
\centering
\caption{Répartition des rôles SCRUM sur le projet}
\label{tab:roles-scrum}
\begin{tabularx}{\textwidth}{L{3.4cm} L{3.4cm} X}
\toprule
\entete{Rôle} & \entete{Tenu par} & \entete{Responsabilité} \\
\midrule
\textit{Product Owner} & Madame Fati, gérante &
Exprime le besoin, arbitre les priorités, valide chaque incrément à l'aune du
fonctionnement réel de l'institut. \\

\textit{SCRUM Master} & \nouvelleEncadrantPedagogique &
Cadence les revues, lève les obstacles et veille au respect de la méthode. \\

\textit{Équipe de développement} & \nouvelleAuteur &
Conception, développement des trois applications, tests, documentation et
déploiement. \\
\bottomrule
\end{tabularx}
\end{table}

\subsection{Réunions}

Trois rendez-vous ont rythmé chaque sprint.

\begin{itemize}[leftmargin=1.2cm]
  \item \textbf{La planification de sprint.} En début de sprint, les éléments
  les plus prioritaires du \textit{product backlog} sont sélectionnés pour
  constituer le \textit{sprint backlog}.

  \item \textbf{La revue de sprint.} En fin de sprint, l'incrément est
  démontré : l'application est installée sur le téléphone de la gérante et le
  back-office ouvert devant elle. C'est là que les écarts entre le besoin
  exprimé et le besoin réel se révèlent.

  \item \textbf{La rétrospective.} Immédiatement après la revue, elle porte sur
  la façon de travailler et non sur le produit. C'est de l'une d'elles qu'est
  née la règle appliquée dans tout le projet : \emph{tout comportement non
  évident doit être expliqué dans le code, à l'endroit où il s'applique}. Cette
  règle explique la densité inhabituelle de commentaires du dépôt, et elle a
  été décisive lors des reprises après interruption.
\end{itemize}

\subsection{Le développement piloté par les tests}

Une approche inspirée du TDD a été appliquée, non pas uniformément, mais là où
elle apporte le plus : les règles métier délicates, et tout particulièrement
les comportements en situation de concurrence.

Écrire d'abord le test qui prouve que trois réservations simultanées ne peuvent
pas occuper deux cabines oblige à formuler la propriété avant de choisir la
technique. C'est ce test, écrit contre une véritable base PostgreSQL, qui a
imposé le recours à un verrou consultatif : la solution intuitive --- compter
puis insérer --- le faisait échouer de façon reproductible.

\figProjet[0.62]{tdd.png}{Le cycle du développement piloté par les tests}{tdd}

Le même cycle a été appliqué au crédit de fidélité, à l'émission des bons de
récompense, à la rotation des jetons de session et à l'anonymisation des
comptes. Ces cinq domaines concentrent l'essentiel des huit suites de tests
qui écrivent réellement en base.

\subsection{Les sprints}

Le projet s'est déroulé sur six sprints, du 6 juillet au 27 août 2026.

\begin{table}[H]
\centering
\caption{Découpage du projet en sprints}
\label{tab:sprints}
\begin{tabularx}{\textwidth}{L{1.5cm} L{2.9cm} X}
\toprule
\entete{Sprint} & \entete{Période} & \entete{Contenu livré} \\
\midrule
1 & 6 -- 9 juillet &
Socle de l'API : authentification JWT, catalogues de prestations et de
produits, module de rendez-vous, horaires d'ouverture, fidélité, récompenses,
gestion des utilisateurs, jetons de rafraîchissement, gestion du fuseau du
centre, documentation Swagger. \\

2 & 12 -- 14 juillet &
Application mobile cliente : catalogues, réservation, fidélité, profil,
rendez-vous. Extraction des composants réutilisables et refonte graphique
complète aux couleurs de l'institut. \\

3 & 16 -- 23 juillet &
Commandes de produits avec paiement à la remise, et back-office complet :
tableau de bord, commandes, rendez-vous, catalogue, fidélité, utilisateurs. \\

4 & 7 -- 11 août &
Durcissement : vérification de l'adresse e-mail, réinitialisation du mot de
passe, récompenses de palier, écritures atomiques sur le solde, le stock et
les créneaux, anonymisation des comptes supprimés, identité native de
l'application mobile. \\

5 & 12 -- 18 août &
Réglages du centre et fermetures exceptionnelles, contrôle du stock,
pagination des listes, finitions d'expérience utilisateur, exports Excel
réservés à la gérante, internationalisation en français, arabe et anglais. \\

6 & 21 -- 27 août &
Hébergement local des images du catalogue, optimisation du décodage des
images, mise à jour du script de sauvegarde, conteneurisation et finalisation
de la documentation d'exploitation. \\
\bottomrule
\end{tabularx}
\end{table}

\subsection{Planification du projet}

La planification a été établie avant le démarrage puis ajustée à chaque revue.
Le diagramme de Gantt ci-dessous en donne la vue d'ensemble : les phases
d'analyse et de conception y précèdent le développement, mais elles ne s'y
substituent pas --- conformément à la démarche agile, une part de conception a
été refaite à chaque sprint, au vu de ce que la revue précédente avait révélé.

\figProjet[0.95]{gantt.png}{Diagramme de Gantt du projet}{gantt}

\section*{Conclusion}
\addcontentsline{toc}{section}{Conclusion}

Ce chapitre a permis de cerner le contexte du projet : un institut de beauté
dont l'activité, entièrement pilotée à la main, souffrait d'une information
unique, localisée et périssable. La solution retenue --- une API portant les
règles métier, une application mobile pour les clientes et un back-office pour
l'institut --- a été justifiée, ses enjeux et ses contraintes énoncés, et la
méthode de conduite du projet exposée. Le chapitre suivant traduit ces besoins
en une analyse et une conception détaillées.

% ═══════════════════════════════════════════════════════════════════
%  Chapitre 2 — Analyse et conception
% ═══════════════════════════════════════════════════════════════════
\chapter{Analyse et conception}

\section*{Introduction}
\addcontentsline{toc}{section}{Introduction}

Ce chapitre présente le résultat du travail d'analyse et de conception. Il
commence par la spécification des besoins --- fonctionnels puis non
fonctionnels --- et par l'identification des acteurs du système. Il expose
ensuite la modélisation UML retenue : les diagrammes de cas d'utilisation,
leur description textuelle, les diagrammes de séquence des processus les plus
délicats, et enfin le diagramme de classes qui fixe la structure des données.

Une précision liminaire sur la méthode : les modèles présentés ici ne sont pas
des artefacts produits en amont puis abandonnés. Ils ont été révisés à chaque
sprint, et le diagramme de classes en particulier reflète l'état final d'un
schéma qui a connu vingt migrations successives. Les choix qui l'ont fait
évoluer sont explicités en fin de chapitre, car ils constituent l'essentiel de
l'apport de conception de ce projet.

% ─────────────────────────────────────────────────────────────────
\section{Spécifications des besoins}

\subsection{Besoins fonctionnels}

Les besoins fonctionnels décrivent ce que le système doit permettre de faire.
Ils ont été regroupés en huit domaines.

\paragraph{Gestion des comptes et de l'authentification.}
\begin{itemize}[leftmargin=1.2cm,nosep]
  \item Inscription d'une cliente avec adresse e-mail, mot de passe, nom,
  prénom et téléphone.
  \item Confirmation de l'adresse e-mail par un code à usage unique.
  \item Connexion, déconnexion et maintien de session sans reconnexion
  quotidienne.
  \item Réinitialisation du mot de passe oublié, par code envoyé par e-mail.
  \item Consultation et modification du profil, changement de mot de passe.
  \item Suppression du compte à la demande de la cliente.
\end{itemize}

\paragraph{Catalogue de prestations.}
\begin{itemize}[leftmargin=1.2cm,nosep]
  \item Consultation publique des prestations actives, par catégorie.
  \item Consultation du détail d'une prestation : description, tarif, photo.
  \item Création, modification, activation et désactivation des prestations et
  de leurs catégories par la gérante.
\end{itemize}

\paragraph{Catalogue de produits.}
\begin{itemize}[leftmargin=1.2cm,nosep]
  \item Consultation publique des produits actifs, par catégorie.
  \item Affichage de la disponibilité en stock.
  \item Gestion complète du catalogue et des stocks par la gérante.
\end{itemize}

\paragraph{Gestion des rendez-vous.}
\begin{itemize}[leftmargin=1.2cm,nosep]
  \item Consultation des créneaux disponibles pour une prestation à une date
  donnée.
  \item Réservation d'un créneau par la cliente.
  \item Consultation de ses propres rendez-vous, à venir et passés.
  \item Annulation par la cliente.
  \item Réservation \emph{pour le compte} d'une cliente par le personnel, pour
  les rendez-vous pris par téléphone ou au comptoir.
  \item Confirmation, clôture, annulation et report d'un rendez-vous par le
  personnel.
\end{itemize}

\paragraph{Gestion des commandes.}
\begin{itemize}[leftmargin=1.2cm,nosep]
  \item Panier et passage de commande, avec retrait à l'institut ou livraison à
  domicile.
  \item Paiement à la remise ou à la livraison.
  \item Consultation de ses commandes et de leur avancement par la cliente.
  \item Annulation par la cliente tant que la commande n'a pas été remise.
  \item Traitement des commandes par le personnel selon un cycle de vie
  contrôlé.
\end{itemize}

\paragraph{Programme de fidélité.}
\begin{itemize}[leftmargin=1.2cm,nosep]
  \item Gain automatique de points sur les prestations réalisées et les
  commandes remises.
  \item Consultation du solde et de l'historique des mouvements.
  \item Catalogue de récompenses échangeables contre des points.
  \item Paliers de visites offrant automatiquement une récompense, avec option
  de récurrence.
  \item Émission d'un bon de récompense doté d'un code lisible, honoré au
  comptoir.
  \item Ajout manuel de points par le personnel, avec motif obligatoire et
  journal consultable par la gérante.
\end{itemize}

\paragraph{Paramétrage du centre.}
\begin{itemize}[leftmargin=1.2cm,nosep]
  \item Définition des horaires d'ouverture, à raison de plusieurs plages par
  jour.
  \item Déclaration de fermetures exceptionnelles (congés, jours fériés).
  \item Réglage du nombre de places réservables simultanément et de l'écart
  entre créneaux, sans redéploiement.
\end{itemize}

\paragraph{Pilotage et exports.}
\begin{itemize}[leftmargin=1.2cm,nosep]
  \item Tableau de bord : commandes en attente, rendez-vous du jour, alertes de
  stock faible.
  \item Gestion des comptes du personnel et attribution des rôles.
  \item Export au format Excel des clientes, des commandes et des rendez-vous,
  réservé à la gérante.
\end{itemize}

\subsection{Besoins non fonctionnels}

Les besoins non fonctionnels portent sur les qualités attendues de la solution
plutôt que sur ses fonctions. Sept ont été retenus, chacun assorti du moyen
concret par lequel il est satisfait --- un besoin non fonctionnel qui ne se
traduit par aucune décision d'implémentation n'est qu'une intention.

\begin{table}[H]
\centering
\caption{Besoins non fonctionnels et moyens de satisfaction}
\label{tab:besoins-non-fonctionnels}
\begin{tabularx}{\textwidth}{L{3.1cm} X}
\toprule
\entete{Besoin} & \entete{Moyen retenu} \\
\midrule
\textbf{Exactitude} &
Aucun créneau ni aucune unité de stock ne peut être attribué deux fois. Les
réservations concurrentes sont sérialisées par un verrou consultatif
PostgreSQL ; le stock et le solde de points sont modifiés par écritures
atomiques conditionnelles. \\

\textbf{Sécurité} &
Mots de passe hachés avec Argon2 ; jetons d'accès JWT de courte durée ; jetons
de rafraîchissement stockés hachés, avec rotation et détection de réutilisation ;
autorisation par rôle vérifiée côté serveur sur chaque route ; en-têtes de
sécurité HTTP ; validation stricte de toute entrée ; limitation du débit par
adresse IP. \\

\textbf{Fiabilité} &
Aucune erreur technique n'atteint l'utilisatrice : un filtre d'exceptions
global traduit les erreurs de base de données en messages compréhensibles.
Les états des rendez-vous et des commandes sont régis par des machines à états
explicites qui interdisent toute transition incohérente. \\

\textbf{Performance} &
Toutes les listes susceptibles de croître sont paginées côté serveur ; le
tableau de bord ne demande que ce qu'il affiche ; les images du catalogue sont
servies avec un cache d'un an et décodées à la taille réellement affichée. \\

\textbf{Ergonomie} &
Charte graphique unifiée entre les trois applications, modes clair et sombre,
messages en langue naturelle, aucun terme technique exposé. \\

\textbf{Accessibilité linguistique} &
Interface mobile disponible en français, en arabe et en anglais ; les réponses
de l'API suivent l'en-tête \texttt{Accept-Language} de la requête et les
e-mails la langue choisie par la cliente. \\

\textbf{Conformité} &
Droit à l'effacement réellement exerçable, par anonymisation du compte plutôt
que par suppression en cascade, afin de préserver l'historique comptable de
l'institut (loi 09-08, CNDP). \\
\bottomrule
\end{tabularx}
\end{table}

\subsection{Acteurs du système}

Cinq acteurs interagissent avec le système : quatre acteurs humains --- la
visiteuse non authentifiée et les trois rôles applicatifs --- et un acteur
secondaire non humain.

\begin{itemize}[leftmargin=1.2cm]
  \item \textbf{La visiteuse} --- acteur non authentifié. Elle consulte les
  catalogues publics et peut créer un compte. Elle ne peut ni réserver ni
  commander.

  \item \textbf{La cliente} (\texttt{CLIENT}) --- acteur principal de
  l'application mobile. Elle réserve, commande, suit son compte de fidélité et
  gère son profil. Elle n'a aucun accès au back-office : une tentative de
  connexion y est refusée à l'entrée.

  \item \textbf{Le personnel} (\texttt{STAFF}) --- acteur principal du
  back-office. Il traite les commandes et les rendez-vous, réserve pour le
  compte d'une cliente, honore les bons de récompense et crédite manuellement
  des points.

  \item \textbf{La gérante} (\texttt{ADMIN}) --- elle dispose de tous les
  droits du personnel, auxquels s'ajoutent la maîtrise du catalogue, des
  horaires, des réglages du centre, du programme de fidélité, des comptes
  utilisateurs, du journal des ajouts manuels de points et des exports Excel.

  \item \textbf{Le service de messagerie} --- acteur secondaire non humain. Il
  achemine les codes de vérification d'adresse et de réinitialisation de mot de
  passe.
\end{itemize}

\subsection{Langage de modélisation}

Le langage UML (\textit{Unified Modeling Language}) a été retenu pour
modéliser le système. C'est un langage graphique normalisé, indépendant de tout
processus de développement et de tout langage de programmation, ce qui permet
de l'intégrer sans friction à une démarche agile : les diagrammes sont révisés
à chaque sprint plutôt que figés au départ.

Trois types de diagrammes ont été employés, chacun pour ce qu'il montre le
mieux :

\begin{itemize}[leftmargin=1.2cm]
  \item les \textbf{diagrammes de cas d'utilisation}, pour délimiter le
  périmètre fonctionnel et répartir les droits entre acteurs ;
  \item les \textbf{diagrammes de séquence}, pour décrire l'ordre des échanges
  dans les processus où cet ordre est précisément ce qui fait la correction du
  système ;
  \item le \textbf{diagramme de classes}, pour fixer la structure des données
  persistantes et leurs relations.
\end{itemize}

% ─────────────────────────────────────────────────────────────────
\section{Diagrammes de cas d'utilisation}

\subsection{Diagramme de cas d'utilisation général}

La figure ci-dessous présente la vue d'ensemble des cas d'utilisation du
système, tous acteurs confondus.

\figProjet{cu-general.png}{Diagramme de cas d'utilisation général de la plate-forme}{cu-general}

\subsection{Cas d'utilisation : réservation d'une prestation}

Le domaine de la réservation est le plus riche du système. Il fait intervenir
la cliente et le personnel sur des cas partiellement communs : le personnel
peut réserver pour le compte d'une cliente, ce qui est le même cas
d'utilisation exercé par un autre acteur, mais il dispose en outre du report et
du changement de statut.

\figProjet{cu-reservation.png}{Diagramme de cas d'utilisation --- gestion des rendez-vous}{cu-reservation}

\subsection{Cas d'utilisation : commande de produits}

\figProjet{cu-commande.png}{Diagramme de cas d'utilisation --- commande de produits}{cu-commande}

\subsection{Cas d'utilisation : programme de fidélité}

Le programme de fidélité comporte deux circuits de récompense bien
distincts, qui convergent vers un même objet : le bon de récompense. Le premier
est automatique : un palier de visites franchi débloque une récompense que la
cliente réclame. Le second est volontaire : la cliente échange ses points
contre une récompense du catalogue. Dans les deux cas, un bon est émis, et
c'est ce bon que le personnel honore au comptoir.

\figProjet{cu-fidelite.png}{Diagramme de cas d'utilisation --- programme de fidélité}{cu-fidelite}

\subsection{Cas d'utilisation : administration}

\figProjet{cu-backoffice.png}{Diagramme de cas d'utilisation --- back-office}{cu-backoffice}

% ─────────────────────────────────────────────────────────────────
\section{Description des cas d'utilisation}

Les cas d'utilisation les plus structurants sont décrits ci-après sous forme
textuelle. Les scénarios alternatifs y occupent une place volontairement
importante : ce sont eux qui déterminent la robustesse du système, et plusieurs
d'entre eux n'ont été identifiés qu'à l'usage.

\begin{table}[H]
\centering
\caption{Description du cas d'utilisation « Réserver une prestation »}
\label{tab:cu-reserver}
\begin{tabularx}{\textwidth}{L{3.1cm} X}
\toprule
\entete{Cas d'utilisation} & Réserver une prestation \\
\midrule
\entete{Acteur principal} & La cliente \\
\entete{Objectif} & Obtenir un rendez-vous pour une prestation, à une date et
une heure disponibles \\
\entete{Préconditions} & La cliente est authentifiée ; la prestation est active
au catalogue \\
\entete{Scénario nominal} &
\begin{enumerate}[leftmargin=*,nosep]
  \item La cliente sélectionne une prestation et une date.
  \item Le système calcule la grille des créneaux du jour à partir des plages
  d'ouverture, de l'écart réglé par le centre et des fermetures
  exceptionnelles.
  \item Le système marque comme indisponible tout créneau dont la capacité est
  atteinte, ainsi que tout créneau déjà passé.
  \item La cliente choisit un créneau libre et valide.
  \item Le système prend le verrou de réservation, contrôle à nouveau la
  disponibilité, enregistre le rendez-vous au statut \textit{en attente} et fige
  le tarif de la prestation.
  \item Le système affiche la confirmation.
\end{enumerate} \\
\entete{Alternatives} &
\textbf{2a.} Le jour est fermé (aucune plage, ou fermeture exceptionnelle) : le
système l'indique et ne propose aucun créneau.\newline
\textbf{4a.} Le créneau a été pris entre l'affichage et la validation : le
système refuse et invite à en choisir un autre.\newline
\textbf{5a.} La date demandée est passée : le système refuse la réservation. \\
\entete{Postconditions} & Le rendez-vous existe au statut \textit{en attente} ;
la capacité du créneau est diminuée d'une unité \\
\bottomrule
\end{tabularx}
\end{table}

\begin{table}[H]
\centering
\caption{Description du cas d'utilisation « Passer une commande »}
\label{tab:cu-commander}
\begin{tabularx}{\textwidth}{L{3.1cm} X}
\toprule
\entete{Cas d'utilisation} & Passer une commande de produits \\
\midrule
\entete{Acteur principal} & La cliente \\
\entete{Objectif} & Commander des produits, avec retrait à l'institut ou
livraison à domicile \\
\entete{Préconditions} & La cliente est authentifiée ; son panier n'est pas
vide \\
\entete{Scénario nominal} &
\begin{enumerate}[leftmargin=*,nosep]
  \item La cliente ajoute des produits à son panier et ouvre la validation.
  \item Elle choisit le mode de remise ; en cas de livraison, elle saisit une
  adresse.
  \item Le système fusionne les lignes en double, contrôle que chaque produit
  est actif et que le stock est suffisant.
  \item Le système fige le prix unitaire de chaque ligne, calcule le total et
  enregistre la commande au statut \textit{en attente}.
  \item La cliente reçoit le récapitulatif de sa commande.
\end{enumerate} \\
\entete{Alternatives} &
\textbf{2a.} Livraison choisie sans adresse : le système refuse.\newline
\textbf{3a.} Un produit a été retiré du catalogue ou son stock est
insuffisant : le système refuse la commande en désignant le produit en
cause. \\
\entete{Postconditions} & La commande existe au statut \textit{en attente} ;
\emph{le stock n'est pas encore décrémenté} --- il ne l'est qu'à la
confirmation par le personnel \\
\bottomrule
\end{tabularx}
\end{table}

Ce dernier point mérite d'être souligné, car il résulte d'un arbitrage
explicite. Décrémenter le stock dès la commande aurait permis de le réserver,
mais aurait ouvert la porte au blocage du stock par des commandes jamais
honorées. Le paiement s'effectuant à la remise, l'institut confirme chaque
commande par un appel : c'est ce moment, et non la commande elle-même, qui
engage réellement le produit.

\begin{table}[H]
\centering
\caption{Description du cas d'utilisation « Échanger des points contre une récompense »}
\label{tab:cu-echanger}
\begin{tabularx}{\textwidth}{L{3.1cm} X}
\toprule
\entete{Cas d'utilisation} & Échanger des points contre une récompense \\
\midrule
\entete{Acteur principal} & La cliente \\
\entete{Objectif} & Convertir un solde de points en récompense utilisable à
l'institut \\
\entete{Préconditions} & La cliente est authentifiée ; la récompense est active
au catalogue \\
\entete{Scénario nominal} &
\begin{enumerate}[leftmargin=*,nosep]
  \item La cliente consulte le catalogue des récompenses et son solde.
  \item Elle demande l'échange d'une récompense.
  \item Le système débite le solde de façon atomique, conditionnellement à sa
  suffisance.
  \item Le système enregistre le mouvement et émet un bon de récompense muni
  d'un code lisible.
  \item La cliente présente ce code au comptoir lors de sa prochaine visite.
\end{enumerate} \\
\entete{Alternatives} &
\textbf{3a.} Le solde est devenu insuffisant entre l'affichage et la demande :
le débit conditionnel n'affecte aucune ligne, l'échange est refusé et rien
n'est émis. \\
\entete{Postconditions} & Le solde est diminué du coût de la récompense ; un
bon dû figure sur la liste du comptoir \\
\bottomrule
\end{tabularx}
\end{table}

\begin{table}[H]
\centering
\caption{Description du cas d'utilisation « Traiter une commande »}
\label{tab:cu-traiter}
\begin{tabularx}{\textwidth}{L{3.1cm} X}
\toprule
\entete{Cas d'utilisation} & Traiter une commande \\
\midrule
\entete{Acteur principal} & Le personnel \\
\entete{Objectif} & Faire progresser une commande jusqu'à sa remise à la
cliente \\
\entete{Préconditions} & Le membre du personnel est authentifié au back-office
avec le rôle \texttt{STAFF} ou \texttt{ADMIN} \\
\entete{Scénario nominal} &
\begin{enumerate}[leftmargin=*,nosep]
  \item Le personnel ouvre la liste des commandes en attente.
  \item Il confirme la commande après appel à la cliente : le système
  décrémente le stock de chaque ligne de façon atomique.
  \item Il marque la commande \textit{prête} lorsqu'elle est préparée.
  \item Il la marque \textit{terminée} à la remise : le système crédite les
  points de fidélité correspondants et met à jour le compteur de visites.
\end{enumerate} \\
\entete{Alternatives} &
\textbf{2a.} Le stock est devenu insuffisant : la confirmation échoue et la
commande reste en attente.\newline
\textbf{$\ast$a.} Une transition non autorisée est demandée (par exemple
\textit{terminée} depuis \textit{en attente}) : le système la refuse.\newline
\textbf{$\ast$b.} La commande est annulée après confirmation : le stock
décrémenté est restitué. \\
\entete{Postconditions} & La commande est au statut \textit{terminée} ; les
points sont crédités exactement une fois \\
\bottomrule
\end{tabularx}
\end{table}

% ─────────────────────────────────────────────────────────────────
\section{Diagrammes de séquence}

Cinq processus ont été modélisés sous forme de diagrammes de séquence. Ils ont
été choisis non pour leur fréquence d'usage, mais parce que dans chacun d'eux
\emph{l'ordre des opérations détermine la correction du résultat} --- ce qu'un
diagramme de cas d'utilisation ne peut pas montrer.

\subsection{Inscription et vérification de l'adresse e-mail}

\figProjet{seq-inscription.png}{Diagramme de séquence --- inscription et vérification de l'adresse}{seq-inscription}

L'inscription crée le compte, son compte de fidélité associé, et renvoie
immédiatement les jetons de session. Ce dernier point n'allait pas de soi : une
première version renvoyait un simple accusé de création, et l'application
enchaînait sur un appel de connexion. Entre les deux appels, la protection
anti-force brute pouvait refuser le second, laissant la cliente devant un
message d'échec alors même que son compte venait d'être créé. L'inscription
renvoie donc désormais exactement ce que renvoie la connexion.

Un code à six chiffres est envoyé par e-mail dans la langue choisie par la
cliente. Seule son empreinte est conservée en base : une fuite de la base ne
livrerait aucun code exploitable. Le nombre de tentatives est plafonné --- six
chiffres représentent un million de possibilités, ce qui se parcourt en
quelques minutes sans ce garde-fou. L'émission d'un nouveau code efface le
précédent du même usage, ce qui garantit qu'au plus une ligne active existe par
cliente et par usage.

\subsection{Réservation d'un créneau}

C'est le processus le plus délicat du système, et le diagramme ci-dessous en
montre la raison.

\figProjet{seq-reservation.png}{Diagramme de séquence --- réservation d'un créneau}{seq-reservation}

La difficulté tient à ce que la contrainte de capacité porte sur un
\emph{ensemble} de lignes --- le nombre de rendez-vous qui chevauchent le
créneau --- et non sur une ligne particulière. Aucune écriture conditionnelle
ne peut donc la garantir. La séquence « compter les rendez-vous du créneau,
puis insérer si la capacité le permet » est correcte pour une cliente seule, et
fausse dès que deux réservations se croisent : les deux comptent avant que
l'une ou l'autre n'ait inséré, les deux trouvent de la place, et deux cabines
en accueillent trois.

La solution retenue consiste à ouvrir une transaction et à y prendre un verrou
consultatif PostgreSQL, tenu jusqu'à la fin de la transaction. La seconde
réservation attend que la première soit écrite avant de compter. Ce verrou se
libère automatiquement au \textit{commit} comme au \textit{rollback}, ce qui
écarte tout risque de blocage résiduel.

À l'intérieur de cette transaction, la vérification enchaîne quatre contrôles :
la prestation existe et est active, la date n'est pas passée, le jour n'est ni
fermé ni exceptionnellement clos, et la capacité du créneau n'est pas atteinte.
Le tarif est enfin figé sur le rendez-vous, afin qu'une révision ultérieure des
prix ne modifie pas les points gagnés sur un rendez-vous déjà réservé.

\subsection{Traitement d'une commande et gestion du stock}

\figProjet{seq-commande.png}{Diagramme de séquence --- commande et décrément du stock}{seq-commande}

Le décrément du stock à la confirmation obéit à la même exigence d'atomicité
que la réservation, mais il se traite plus simplement : la contrainte porte
cette fois sur une ligne unique --- le produit --- et une mise à jour
conditionnelle suffit. La condition « décrémenter de $n$ à condition que le
stock soit au moins égal à $n$ » est évaluée par la base au moment de
l'écriture. Si elle n'est pas satisfaite, aucune ligne n'est affectée et la
transaction entière est annulée : la commande reste en attente et le stock
demeure cohérent.

\subsection{Clôture d'un rendez-vous et crédit de fidélité}

\figProjet{seq-fidelite.png}{Diagramme de séquence --- clôture d'un rendez-vous, crédit de points et palier}{seq-fidelite}

La clôture d'un rendez-vous déclenche une chaîne d'opérations qui doivent
toutes réussir ou toutes échouer : passage au statut \textit{terminé}, calcul
des points sur le tarif figé, incrément du solde et du compteur de visites,
écriture du mouvement d'historique, puis évaluation des paliers de visites. Le
tout s'exécute dans une seule transaction.

L'idempotence est assurée à deux niveaux. Le mouvement de fidélité porte une
référence unique vers le rendez-vous : une clôture rejouée ne peut pas créditer
deux fois. Les déblocages de paliers portent une contrainte d'unicité sur le
triplet (compte, palier, cycle), où le cycle indique combien de fois le seuil a
été franchi --- ce qui permet à un palier d'être récurrent tout en gardant
chaque déblocage unique.

\subsection{Rotation des jetons de session}

\figProjet{seq-refresh.png}{Diagramme de séquence --- rafraîchissement et rotation des jetons}{seq-refresh}

Le jeton d'accès est volontairement de courte durée. Pour éviter de demander à
la cliente de se reconnecter plusieurs fois par jour, un jeton de
rafraîchissement de longue durée permet d'en obtenir un nouveau.

Ce mécanisme est protégé par une rotation : chaque rafraîchissement révoque le
jeton présenté et en émet un nouveau, rattaché à la même famille. Si un jeton
déjà utilisé est présenté une seconde fois, c'est le signe qu'il a été
dérobé --- soit par l'attaquant, soit par la cliente légitime, sans qu'il soit
possible de savoir lequel. Le système révoque alors toute la famille : les deux
parties sont déconnectées, et l'accès frauduleux prend fin.

% ─────────────────────────────────────────────────────────────────
\section{Diagramme de classes}

\subsection{Vue d'ensemble}

Le diagramme de classes ci-dessous présente la structure des données
persistantes. Il comporte dix-neuf entités, réparties en six domaines :
utilisateurs et sessions, catalogues, rendez-vous, commandes, fidélité, et
paramétrage du centre.

\figProjet[0.98]{diagramme-classes.png}{Diagramme de classes de la plate-forme}{diagramme-classes}

\subsection{Choix structurants du modèle}

Cinq décisions de modélisation méritent d'être explicitées, car elles ne
découlent pas mécaniquement des besoins et constituent l'essentiel de l'apport
de conception.

\paragraph{Le figement des montants.} Un rendez-vous conserve le tarif de la
prestation \emph{tel qu'il était à la réservation}, et chaque ligne de commande
conserve le prix unitaire du produit \emph{tel qu'il était à la commande}.
Sans cela, une révision des tarifs réécrirait rétroactivement l'historique :
une commande passée en juillet afficherait le prix d'août, et les points gagnés
sur un rendez-vous changeraient après coup. Le même principe s'applique aux
récompenses de palier, dont la référence est figée au déblocage, et aux points
dépensés, figés sur le bon émis.

\paragraph{L'anonymisation plutôt que la suppression.} Une cliente qui exerce
son droit à l'effacement voit son compte vidé de tout ce qui l'identifie ---
adresse e-mail, nom, téléphone, mot de passe --- et marqué comme supprimé, mais
sa ligne n'est pas détruite. Ses commandes et ses rendez-vous portent en effet
le chiffre d'affaires de l'institut, qui doit survivre à son départ. Les
relations le garantissent structurellement : les rendez-vous et les commandes
sont rattachés à la cliente en mode \texttt{Restrict}, si bien qu'une
suppression échouerait bruyamment plutôt que de détruire l'historique en
silence. Un compte ainsi marqué ne peut plus se connecter, ni recevoir
d'e-mail, ni réinitialiser son mot de passe.

\paragraph{Le traitement différencié des employées.} Les écritures de fidélité
saisies par le personnel et les bons honorés au comptoir désignent l'employée
concernée, mais en mode \texttt{SetNull} et non \texttt{Cascade}. Le départ
d'une employée ne doit pas effacer les points des clientes qu'elle a servies ni
la trace de ce qui leur a été remis : l'écriture demeure, elle ne désigne
simplement plus personne.

\paragraph{Le bon de récompense comme entité de première classe.} C'est la
notion qui manquait à la première version du programme de fidélité. Sans elle,
réclamer une récompense offerte la faisait disparaître des deux côtés à la fois
--- l'application filtrait sur la date de réclamation, la fiche du comptoir
aussi --- et un échange de points ne laissait qu'une ligne d'historique que
personne ne listait. Le bon réunit les deux circuits de récompense en un objet
unique : une seule chose à afficher pour la cliente, une seule liste à traiter
au comptoir. Son unicité par origine --- un déblocage de palier, ou un
mouvement de points --- constitue le véritable garde-fou d'idempotence : une
réclamation rejouée ne peut pas émettre un second bon.

\paragraph{La plage d'ouverture plutôt que le jour ouvré.} Les horaires sont
modélisés comme des \emph{plages}, et non comme des jours dotés d'une heure
d'ouverture et d'une heure de fermeture. Un même jour en porte donc autant que
nécessaire --- 9\,h--13\,h puis 14\,h--18\,h --- et un jour sans aucune plage
est fermé. Ce modèle exprime naturellement la pause déjeuner, qu'un booléen
« fermé » assorti de deux heures ne savait représenter qu'au prix d'une plage
trouée. Les fermetures exceptionnelles, elles, sont stockées comme des dates de
calendrier et non comme des instants : une fermeture est un jour du calendrier
local du centre, et stocker un instant obligerait à choisir une heure
arbitraire.

\section*{Conclusion}
\addcontentsline{toc}{section}{Conclusion}

Ce chapitre a présenté l'analyse et la conception du système : les besoins
fonctionnels et non fonctionnels, les acteurs, les cas d'utilisation et leur
description, les séquences des processus les plus délicats, et la structure des
données. Les décisions de modélisation y ont été justifiées plutôt que
seulement énoncées, car plusieurs d'entre elles --- le figement des montants,
l'anonymisation, le bon de récompense --- constituent la réponse à des
difficultés qui ne se voyaient qu'à l'usage. Le chapitre suivant expose les
choix techniques qui permettent de réaliser cette conception.

% ═══════════════════════════════════════════════════════════════════
%  Chapitre 3 — Étude technique
% ═══════════════════════════════════════════════════════════════════
\chapter{Étude technique}

\section*{Introduction}
\addcontentsline{toc}{section}{Introduction}

Ce chapitre expose les choix techniques du projet. Il présente les langages,
les cadriciels et le système de gestion de base de données retenus, ainsi que
les environnements matériel et logiciel de travail. Il décrit ensuite
l'architecture technique de la solution : son découpage en couches, son
architecture de déploiement, et la chaîne de sécurité appliquée à chaque
requête.

Les choix sont ici justifiés plutôt qu'énumérés. Une technologie ne vaut que
par le problème qu'elle résout dans un contexte donné, et plusieurs des options
retenues n'auraient pas été les mêmes pour un projet d'une autre nature.

% ─────────────────────────────────────────────────────────────────
\section{Technologies de développement}

\subsection{Langages}

\paragraph{TypeScript.} L'ensemble du code applicatif --- API, back-office et
application mobile --- est écrit en TypeScript, sur-ensemble typé de
JavaScript développé par Microsoft. Ce choix unique pour les trois applications
a une conséquence directe et précieuse : les types décrivant les données
métier sont écrits une fois et compris partout. Le type d'un rendez-vous
retourné par l'API est le même côté serveur, côté back-office et côté
application mobile.

Le typage statique joue par ailleurs un rôle de filet de sécurité dans un
projet qui n'a pas de tests d'interface : le back-office et l'application
mobile n'ont pas de suite de tests automatisés, et leur garantie principale est
la compilation. Un champ renommé côté serveur produit une erreur de compilation
côté client, avant tout déploiement.

\paragraph{SQL.} Le langage SQL n'est presque jamais écrit à la main dans le
projet --- Prisma le génère --- à une exception près, précisément là où il est
irremplaçable : la prise du verrou consultatif qui sérialise les réservations
concurrentes, exprimée par un appel direct à \texttt{pg\_advisory\_xact\_lock}.
Aucune abstraction ne fournit cette primitive, et c'est elle qui garantit
qu'un créneau n'est jamais vendu deux fois.

\subsection{Cadriciels et bibliothèques}

\paragraph{NestJS.} L'API repose sur NestJS, cadriciel Node.js structuré autour
de modules, de contrôleurs et de services, et bâti sur l'injection de
dépendances. Sa structure imposée est un atout dans un projet mené seul : elle
évite les choix d'organisation arbitraires et rend chaque module immédiatement
localisable. Ses gardes, filtres, intercepteurs et tubes de validation
permettent en outre de traiter les préoccupations transversales --- droits
d'accès, validation, traduction, gestion des erreurs --- à un seul endroit,
plutôt que répétées dans chaque contrôleur.

\paragraph{Prisma.} Prisma assure la correspondance objet-relationnel. Le
schéma y est déclaré dans un fichier unique, à partir duquel sont générés à la
fois le client typé et les fichiers de migration. Deux propriétés ont pesé dans
ce choix : le client est \emph{typé d'après le schéma}, si bien qu'une colonne
renommée fait échouer la compilation partout où elle était lue ; et les
migrations sont des fichiers SQL versionnés, relus avant application, ce qui
rend le passage en production vérifiable.

\paragraph{React et Vite.} Le back-office est une application web mono-page
écrite en React et servie par Vite. React a été retenu pour la richesse de son
écosystème et sa maîtrise préalable ; Vite pour la rapidité de son serveur de
développement et la simplicité de sa production statique. L'état global ---
réduit à la session authentifiée --- est géré par Zustand, bibliothèque
volontairement minimale, dimensionnée pour ce besoin.

\paragraph{React Native et Expo.} L'application mobile est développée avec
React Native, sous la chaîne d'outils Expo. Ce choix permet de produire une
application Android et iOS à partir d'une base de code unique, tout en
partageant le langage et une partie des conventions avec le back-office. Expo
apporte en outre les briques natives nécessaires --- stockage sécurisé des
jetons, polices, images, dégradés, localisation --- sans avoir à écrire de code
natif.

\paragraph{Bibliothèques notables.} Argon2 pour le hachage des mots de passe,
choisi pour sa résistance aux attaques matérielles ; \texttt{date-fns-tz} pour
les conversions entre heure locale du centre et temps universel ; ExcelJS pour
la génération des exports ; Nodemailer pour l'acheminement des e-mails ; Helmet
pour les en-têtes de sécurité HTTP ; \texttt{class-validator} pour la
validation déclarative des entrées ; i18next pour l'internationalisation de
l'application mobile.

\subsection{Services web}

L'API suit le style d'architecture REST, avec JSON pour format d'échange. Ce
choix s'imposait pour plusieurs raisons convergentes :

\begin{itemize}[leftmargin=1.2cm]
  \item \textbf{L'universalité de HTTP.} Deux clients de natures très
  différentes --- une application native et une application web --- consomment
  la même API sans adaptation.

  \item \textbf{L'absence d'état du protocole.} Chaque requête porte son jeton
  d'authentification ; aucun état de session n'est conservé côté serveur, ce
  qui simplifie considérablement le déploiement et la montée en charge.

  \item \textbf{La légèreté du format.} JSON convient particulièrement à un
  client mobile, dont la bande passante et l'autonomie sont des ressources
  comptées.

  \item \textbf{L'outillage.} La documentation OpenAPI est générée directement
  depuis les décorateurs du code, ce qui garantit qu'elle ne diverge pas de
  l'implémentation.
\end{itemize}

L'API compte plus de quatre-vingts routes réparties en douze contrôleurs. Le
tableau ci-dessous en donne la répartition.

\begin{table}[H]
\centering
\caption{Répartition des routes de l'API par domaine}
\label{tab:routes-api}
\begin{tabularx}{\textwidth}{L{3.1cm} C{1.7cm} X}
\toprule
\entete{Domaine} & \entete{Routes} & \entete{Principales opérations} \\
\midrule
\texttt{/auth} & 9 &
Inscription, connexion, profil courant, vérification d'adresse, renvoi de code,
mot de passe oublié, réinitialisation, rafraîchissement, déconnexion \\

\texttt{/users} & 6 &
Profil, modification, changement de mot de passe, suppression de compte, liste
des utilisateurs, attribution de rôle \\

\texttt{/services} & 10 &
Catalogue public, catégories, détail, gestion complète (administration) \\

\texttt{/products} & 10 &
Catalogue public, catégories, détail, gestion complète (administration) \\

\texttt{/appointments} & 7 &
Disponibilités, réservation, mes rendez-vous, annulation, réservation pour une
cliente, changement de statut, report \\

\texttt{/orders} & 6 &
Commande, mes commandes, détail, annulation, liste (personnel), changement de
statut \\

\texttt{/loyalty} & 23 &
Solde, historique, récompenses, échange, paliers, déblocages, réclamation,
bons, gestion (administration), ajout manuel, journal d'audit \\

\texttt{/opening-\allowbreak hours} & 5 &
Consultation, remplacement des plages, fermetures exceptionnelles \\

\texttt{/settings} & 2 & Consultation et modification des réglages du centre \\
\texttt{/exports} & 3 & Exports Excel des clientes, commandes et rendez-vous \\
\texttt{/uploads} & 1 & Téléversement d'une image de catalogue \\
\texttt{/health} & 1 & Sonde de santé pour la supervision et le déploiement \\
\bottomrule
\end{tabularx}
\end{table}

\subsection{Système de gestion de base de données}

\textbf{PostgreSQL 16} a été retenu comme système de gestion de base de
données. Le choix d'un SGBD relationnel ne faisait pas débat --- les données du
projet sont fortement structurées et fortement liées --- mais celui de
PostgreSQL plutôt que d'une autre base relationnelle mérite d'être justifié.

Trois de ses propriétés sont directement exploitées :

\begin{itemize}[leftmargin=1.2cm]
  \item \textbf{Les verrous consultatifs.} C'est la primitive qui rend correcte
  la réservation concurrente. Elle n'a pas d'équivalent direct et portable
  ailleurs, et elle est utilisée au c\oe ur même du moteur de réservation.

  \item \textbf{Les transactions et l'intégrité référentielle.} Le crédit de
  fidélité, la clôture d'un rendez-vous et le déblocage d'un palier
  s'exécutent dans une transaction unique. Les modes de suppression
  différenciés --- \texttt{Cascade}, \texttt{Restrict}, \texttt{SetNull} selon
  la relation --- portent une partie des règles de conservation de
  l'historique dans le schéma lui-même, où aucun oubli applicatif ne peut les
  contourner.

  \item \textbf{Le type décimal exact.} Les montants sont stockés en
  \texttt{Decimal(10,2)}, et non en nombre flottant. Un montant en virgule
  flottante accumule des erreurs d'arrondi qui finissent par se voir sur un
  total.
\end{itemize}

Le schéma comporte dix-neuf tables et a connu vingt migrations versionnées
depuis l'initialisation du projet.

\subsection{Outils de développement}

\subsubsection{Environnement matériel}

\begin{table}[H]
\centering
\caption{Environnement matériel de développement}
\label{tab:env-materiel}
\begin{tabularx}{\textwidth}{L{4.2cm} X}
\toprule
\entete{Élément} & \entete{Caractéristiques} \\
\midrule
Poste de développement & Ordinateur portable, processeur multic\oe ur, 16 Go
de mémoire vive, disque SSD \\
Système d'exploitation & Windows avec sous-système Linux, et Linux \\
Terminal de test mobile & Téléphone Android physique et émulateur Android \\
Serveur de base de données & Conteneur Docker PostgreSQL 16, exécuté localement \\
\bottomrule
\end{tabularx}
\end{table}

\begin{attention}[Un piège rencontré]
Sous Windows, la plage de ports 2933--3032 est réservée par Hyper-V et par le
sous-système Linux. Le port 3000, choisi par défaut par la plupart des projets
Node.js, y produit une erreur d'accès refusé qui ne désigne pas sa cause. Le
port d'écoute de l'API a donc été rendu configurable par variable
d'environnement --- ce qui s'est révélé nécessaire également au déploiement,
la plupart des hébergeurs imposant le leur.
\end{attention}

\subsubsection{Environnement logiciel}

\begin{table}[H]
\centering
\caption{Environnement logiciel de développement}
\label{tab:env-logiciel}
\begin{tabularx}{\textwidth}{L{4.2cm} X}
\toprule
\entete{Outil} & \entete{Usage} \\
\midrule
Visual Studio Code & Éditeur de code des trois applications \\
Node.js 22 & Environnement d'exécution JavaScript côté serveur et outillage \\
Docker et Docker Compose & Base de données locale, puis conteneurisation de
l'API pour le déploiement \\
Git et GitHub & Gestion de versions et hébergement du dépôt \\
GitHub Actions & Intégration continue : compilation et tests à chaque
modification \\
Prisma Studio & Inspection et correction ponctuelle des données en
développement \\
Swagger UI & Documentation interactive de l'API, générée depuis le code \\
Jest et Supertest & Tests unitaires, tests d'intégration et tests de
concurrence \\
ESLint et Prettier & Analyse statique et mise en forme homogène du code \\
Expo Go & Exécution de l'application mobile sur un téléphone réel sans
compilation native \\
Lucidchart & Réalisation des diagrammes UML \\
\bottomrule
\end{tabularx}
\end{table}

% ─────────────────────────────────────────────────────────────────
\section{Architecture technique}

\subsection{Architecture en couches}

Une vue d'ensemble permet de dégager les couches du système. L'application
étant à la fois web et mobile, l'architecture retenue est une architecture en
quatre niveaux : deux clients distincts, un serveur d'application et un serveur
de données.

\figProjet[0.8]{couches.png}{Les couches de l'application}{couches}

\begin{itemize}[leftmargin=1.2cm]
  \item \textbf{La couche de présentation.} Elle encapsule l'interface
  utilisateur. Elle est double : l'application mobile pour les clientes, le
  back-office web pour l'institut. Elle n'implémente \emph{aucune} règle
  métier : elle affiche ce que le serveur lui transmet et lui soumet ce que
  l'utilisatrice saisit. Les contrôles qu'elle effectue --- un champ
  obligatoire, un format d'adresse --- servent uniquement le confort de saisie
  et sont systématiquement rejoués côté serveur.

  \item \textbf{La couche métier.} C'est le serveur d'application NestJS. Elle
  porte l'intégralité des règles : moteur de créneaux, machines à états,
  calcul des points, contrôle des droits, validation des entrées. Elle expose
  ses services sous forme d'API REST.

  \item \textbf{La couche d'accès aux données.} Assurée par Prisma, elle
  traduit les opérations métier en requêtes SQL, gère les transactions et
  garantit le typage des résultats.

  \item \textbf{La couche de persistance.} La base PostgreSQL. Elle ne se
  contente pas de stocker : elle porte des règles d'intégrité --- unicité,
  clés étrangères, modes de suppression --- et fournit les primitives de
  verrouillage qui rendent la concurrence correcte.
\end{itemize}

\subsection{Architecture de déploiement}

\figProjet[0.9]{architecture.png}{Architecture technique et déploiement de la solution}{architecture}

Le déploiement repose sur Docker Compose et enchaîne trois services dans un
ordre imposé :

\begin{enumerate}[leftmargin=1.2cm]
  \item \textbf{PostgreSQL}, doté d'un contrôle de santé et d'un volume
  persistant. Son port n'est publié que sur l'interface locale.

  \item \textbf{Un service de migration}, à usage unique : il applique les
  migrations en attente, puis s'arrête. Il est délibérément séparé de l'API,
  afin qu'une migration en échec se présente comme une étape ratée assortie de
  son message, et non comme une application qui redémarre en boucle sans
  raison apparente. Il utilise la commande \texttt{migrate deploy}, qui
  applique ce qui manque et refuse de toucher à l'existant --- et jamais
  \texttt{migrate dev}, qui peut réinitialiser la base.

  \item \textbf{L'API}, qui ne démarre qu'une fois les migrations appliquées
  avec succès.
\end{enumerate}

L'image de l'API est construite en trois étapes, afin que l'image livrée ne
contienne que ce qui sert à servir --- ni compilateur, ni cadriciel de test,
ni outillage de développement. Le résultat mesuré est de 178 Mo.

\begin{attention}[Une décision de déploiement]
L'API est publiée sur l'interface locale uniquement. C'est délibéré : un
reverse proxy (Caddy ou nginx) doit se placer devant elle pour assurer le
chiffrement HTTPS. Publier le port directement reviendrait à servir l'API en
clair, jetons d'authentification compris. Le jour où ce proxy est mis en place,
un réglage devient indispensable : la limitation de débit compte par adresse
IP, et derrière un proxy toutes les clientes partagent la sienne. Sans ce
réglage, l'institut entier se retrouve sur un compteur unique et l'API commence
à refuser des requêtes au hasard dès qu'il y a du monde --- un symptôme
déroutant, parfait en test avec une seule personne, et que rien dans les
journaux ne rattache à sa cause.
\end{attention}

\subsection{La chaîne de sécurité d'une requête}

La sécurité n'est pas traitée comme une fonctionnalité, mais comme une chaîne
de filtres que toute requête traverse avant d'atteindre la moindre ligne de
logique métier.

\figProjet[0.85]{securite.png}{Chaîne de traitement de sécurité d'une requête}{securite}

\begin{enumerate}[leftmargin=1.2cm]
  \item \textbf{En-têtes de sécurité.} Helmet retire l'en-tête annonçant la
  technologie du serveur et ajoute une douzaine d'en-têtes appliqués par le
  navigateur : transport strict, anti-\textit{sniffing}, anti-\textit{clickjacking},
  politique de sécurité du contenu.

  \item \textbf{Limitation du débit.} Un compteur par adresse IP plafonne
  l'ensemble de l'API à 120 requêtes par minute. Les routes coûteuses portent
  en plus une limite stricte : dix tentatives par minute pour la connexion et
  l'inscription, cinq pour le renvoi de code, et cinq par tranche de cinq
  minutes pour le mot de passe oublié.

  \item \textbf{Contrôle d'origine.} La politique CORS n'autorise que les
  adresses déclarées du back-office.

  \item \textbf{Authentification.} Un garde vérifie la signature et la validité
  du jeton d'accès, sauf sur les routes explicitement publiques.

  \item \textbf{Autorisation.} Un second garde compare le rôle porté par le
  jeton aux rôles exigés par la route. Ce contrôle est \emph{indépendant} de ce
  que l'interface affiche ou masque : une route d'administration atteinte par
  un autre moyen est refusée par le serveur.

  \item \textbf{Validation des entrées.} Un tube de validation global rejette
  toute propriété non déclarée dans le contrat de la route, plutôt que de
  l'ignorer silencieusement, et convertit les valeurs dans leur type attendu.

  \item \textbf{Traitement uniforme des erreurs.} Un filtre global intercepte
  toute exception avant qu'elle n'atteigne le client. Sans lui, les erreurs de
  base de données partaient brutes, avec les noms de tables et de contraintes.
\end{enumerate}

\paragraph{Le stockage des secrets.} Les mots de passe sont hachés avec Argon2,
lauréat de la compétition internationale de hachage de mots de passe. Les
jetons de rafraîchissement et les codes de vérification ne sont jamais stockés
en clair : seule leur empreinte l'est. Une lecture de la base ne livrerait donc
ni mot de passe, ni session utilisable, ni code de réinitialisation.

\paragraph{La défense contre l'énumération par mesure du temps.} Une
vérification de mot de passe avec Argon2 prend un temps mesurable. Si le
serveur répondait immédiatement pour une adresse inconnue et après le calcul
pour une adresse connue, la seule mesure du délai de réponse permettrait de
déterminer quelles adresses possèdent un compte. La connexion exécute donc une
vérification factice sur une valeur aléatoire lorsque l'adresse est inconnue :
les deux cas répondent au même rythme.

\paragraph{Le contrôle des fichiers téléversés.} Les images du catalogue sont
identifiées par leur \emph{signature binaire}, et non par le type déclaré par
le navigateur ni par l'extension du nom de fichier --- deux valeurs qui
viennent du client et se falsifient en une ligne. Trois formats sont acceptés :
JPEG, PNG et WebP. Le SVG en est délibérément absent : c'est du XML susceptible
de contenir du JavaScript, qui s'exécuterait avec les droits du domaine servant
le fichier --- une faille de type XSS ouverte par un simple téléversement.

\paragraph{La validation de la configuration au démarrage.} L'API refuse de
démarrer si sa configuration est incohérente. Le secret de signature des jetons
doit compter au moins trente-deux caractères ; un secret plus court se casse
hors ligne et permet de forger un jeton d'administration sans jamais toucher à
l'API. De même, l'envoi réel d'e-mails est exigé en production : un
environnement de production configuré en mode console n'enverrait aucun code de
vérification, sans qu'aucune erreur ne le signale.

\section*{Conclusion}
\addcontentsline{toc}{section}{Conclusion}

Ce chapitre a présenté les choix techniques du projet et leur justification :
un langage unique pour les trois applications, un cadriciel structurant pour
l'API, un ORM typé adossé à PostgreSQL, et une architecture en quatre niveaux
où la couche métier est seule dépositaire des règles. La chaîne de sécurité
appliquée à chaque requête a été détaillée, ainsi que l'architecture de
déploiement. Le chapitre suivant présente la mise en \oe uvre effective de ces
choix, les tests qui la valident et les interfaces réalisées.

% ═══════════════════════════════════════════════════════════════════
%  Chapitre 4 — Sécurité de la plate-forme
% ═══════════════════════════════════════════════════════════════════
\chapter{Sécurité de la plate-forme}

\section*{Introduction}
\addcontentsline{toc}{section}{Introduction}

Ce chapitre est consacré à la sécurité de la solution. Il occupe une place
particulière dans ce rapport, et pour deux raisons.

La première tient à la nature du système. La plate-forme traite des données
personnelles --- identité, coordonnées, historique de fréquentation d'un
institut de beauté --- et porte des opérations à valeur économique :
réservations, commandes, points de fidélité convertibles en prestations
gratuites. Chacune de ces trois catégories appelle des garanties différentes,
et aucune n'est satisfaite par défaut.

La seconde tient au fait qu'un système exposé sur Internet n'a pas seulement
des utilisatrices : il a des visiteurs. L'API est publique par construction ---
l'application mobile l'appelle depuis les téléphones de clientes, sur des
réseaux quelconques. Tout ce qui n'est pas explicitement refusé côté serveur
est, en pratique, autorisé.

La démarche suivie n'a donc pas consisté à ajouter de la sécurité à un système
achevé, mais à traiter chaque décision de conception comme portant une part de
sécurité. Ce chapitre restitue cette démarche : le modèle de menaces retenu,
les principes directeurs, puis les contre-mesures effectivement implémentées,
domaine par domaine. Il se termine par une synthèse et par un inventaire
honnête des limites qui subsistent.

% ─────────────────────────────────────────────────────────────────
\section{Démarche et modèle de menaces}

\subsection{Surface d'attaque}

La première étape a consisté à délimiter la surface exposée et à situer les
frontières de confiance du système.

\figProjet[0.9]{securite-surface.png}{Surface d'attaque et frontières de confiance}{securite-surface}

Cinq points d'entrée existent, et un seul est réellement de confiance :

\begin{itemize}[leftmargin=1.2cm]
  \item \textbf{L'API HTTP}, atteignable par quiconque connaît son adresse. Sa
  documentation OpenAPI décrit d'ailleurs chacune de ses routes --- raison pour
  laquelle elle est coupée en production.

  \item \textbf{L'application mobile}, qui s'exécute sur le téléphone de la
  cliente. Elle échappe totalement à notre contrôle : son code est lisible, son
  trafic interceptable, ses contrôles contournables. \emph{Aucune règle de
  sécurité ne peut y résider.}

  \item \textbf{Le back-office web}, dans le même cas : c'est du JavaScript
  servi au navigateur. Masquer un bouton d'administration n'a jamais protégé la
  route correspondante.

  \item \textbf{Les images téléversées}, seuls fichiers que l'application écrit
  sur le disque du serveur --- et qu'elle republie ensuite en HTTP.

  \item \textbf{La base de données}, qui n'est pas exposée : son port n'est
  publié que sur l'interface locale du serveur.
\end{itemize}

Il en découle un principe qui structure tout ce chapitre :
\emph{la frontière de confiance passe à l'entrée de l'API, et nulle part
ailleurs}. Tout ce qui arrive d'un client --- corps de requête, en-têtes, type
MIME déclaré, extension de fichier, rôle affiché dans l'interface --- est une
donnée hostile jusqu'à vérification par le serveur.

\subsection{Modèle de menaces}

Les menaces ont été recensées selon les six catégories du modèle STRIDE, en
retenant pour chacune les scénarios réalistes pour ce système.

\begin{table}[H]
\centering
\caption{Modèle de menaces STRIDE appliqué à la plate-forme}
\label{tab:stride}
\begin{tabularx}{\textwidth}{L{2.9cm} X}
\toprule
\entete{Catégorie} & \entete{Scénarios retenus} \\
\midrule
\textbf{Usurpation}\newline\small(\textit{Spoofing}) &
Force brute sur le formulaire de connexion ; vol d'un jeton de session ;
devinette d'un code de réinitialisation à six chiffres ; forge d'un jeton
d'administration si le secret de signature est faible. \\

\textbf{Altération}\newline\small(\textit{Tampering}) &
Injection SQL ; modification du rôle porté dans la requête ; téléversement
d'un fichier exécutable déguisé en image ; falsification de l'en-tête
d'adresse IP réelle. \\

\textbf{Répudiation} &
Ajout manuel de points de fidélité sans trace ; remise d'un bon de récompense
sans enregistrement de qui l'a honoré. \\

\textbf{Divulgation}\newline\small(\textit{Information disclosure}) &
Lecture des commandes ou des rendez-vous d'une autre cliente ; fuite du schéma
de base par les messages d'erreur ; énumération des comptes existants ; export
massif de la base clientes. \\

\textbf{Déni de service} &
Saturation de l'API par un client unique ; épuisement du processeur par le
hachage de mots de passe ; export d'un volume de lignes qui bloque le
serveur ; réservation en masse de créneaux. \\

\textbf{Élévation de privilège} &
Accès direct à une route d'administration depuis un compte cliente ;
promotion de son propre rôle ; usage d'un jeton émis avant la suppression du
compte. \\
\bottomrule
\end{tabularx}
\end{table}

\subsection{Principes directeurs}

Quatre principes ont guidé les décisions.

\paragraph{Le serveur fait autorité.} Toute règle de sécurité est appliquée par
l'API. Les interfaces peuvent masquer, désactiver ou pré-valider : ce n'est
jamais qu'un confort d'usage, jamais une protection.

\paragraph{Défense en profondeur.} Aucune contre-mesure unique n'est réputée
suffisante. Un compte supprimé est ainsi protégé par trois verrous
indépendants : le refus de connexion sur son marqueur de suppression, le
remplacement de son mot de passe par le haché d'un secret que personne ne
connaît, et le rejet, à chaque requête, des jetons déjà émis.

\paragraph{Sécurité par défaut, et échec fermé.} Toute route est authentifiée
et comptée sauf déclaration contraire explicite. Toute propriété non déclarée
dans le contrat d'une route fait rejeter la requête, au lieu d'être ignorée en
silence. Et lorsqu'une configuration est incohérente, l'API \emph{refuse de
démarrer} plutôt que de fonctionner dans un état dégradé que rien ne
signalerait.

\paragraph{Moindre privilège.} Trois rôles seulement, chacun strictement borné.
L'image de production ne contient aucun compilateur et le processus n'y tourne
pas en tant que \texttt{root}. Le port de la base de données n'est pas publié.

% ─────────────────────────────────────────────────────────────────
\section{Authentification}

\subsection{Stockage des mots de passe}

Les mots de passe sont hachés avec \textbf{Argon2id}, lauréat du
\textit{Password Hashing Competition} et normalisé par la RFC 9106. Le choix
s'est porté sur lui plutôt que sur bcrypt pour sa résistance aux attaques
matérielles : Argon2 est conçu pour être coûteux en \emph{mémoire} autant qu'en
temps de calcul, ce qui ôte aux cartes graphiques et aux circuits dédiés
l'avantage massif dont ils disposent face aux fonctions purement calculatoires.

Aucun mot de passe n'est stocké, journalisé ou renvoyé en clair, à aucun
moment. La politique appliquée impose au minimum huit caractères, dont une
majuscule et une minuscule, et elle est \emph{la même} à l'inscription, au
changement de mot de passe et à la réinitialisation --- une exigence relâchée
sur l'un de ces trois chemins annulerait les deux autres.

\subsection{Jetons d'accès}

L'authentification des requêtes repose sur des jetons JWT signés, de durée
volontairement courte : \textbf{quinze minutes}. Cette brièveté est la réponse
au fait qu'un JWT est, par nature, invérifiable auprès d'une liste de
révocation sans requête supplémentaire : elle borne la fenêtre pendant laquelle
un jeton dérobé reste exploitable.

Le secret de signature est lu par une méthode qui \emph{lève une exception s'il
est absent}, aussi bien à la signature qu'à la vérification. C'est délibéré :
une valeur de repli en dur --- le réflexe habituel pour « que ça démarre en
développement » --- rendrait tous les jetons forgeables par quiconque a lu le
dépôt.

Le jeton n'est cependant pas cru sur parole. À chaque requête, la stratégie de
vérification recharge le compte en base et refuse le jeton si le compte
n'existe plus ou a été anonymisé.

\begin{lstlisting}[caption={Vérification du jeton d'accès : la signature ne suffit pas},captionpos=b]
async validate(payload: { sub: string; email: string; role: string }) {
  const user = await this.usersService.findById(payload.sub);
  // Un compte anonymise est traite comme inexistant. Sans ce test, l'access
  // token emis juste avant la suppression continuerait d'ouvrir le compte
  // jusqu'a son expiration.
  if (!user || user.deletedAt) {
    throw new UnauthorizedException();
  }
  return { id: user.id, email: user.email, role: user.role };
}
\end{lstlisting}

Ce rechargement a un second effet, moins évident et tout aussi important : le
rôle appliqué est celui \emph{de la base}, et non celui inscrit dans le jeton.
Une rétrogradation prend donc effet immédiatement, sans attendre l'expiration
des jetons en circulation.

\subsection{Jetons de rafraîchissement, rotation et détection de vol}

Des jetons d'accès de quinze minutes obligeraient la cliente à se reconnecter
plusieurs fois par jour. Un jeton de rafraîchissement de longue durée résout ce
problème --- et en crée un autre : c'est désormais lui qui vaut la session, et
sa durée de vie le rend d'autant plus intéressant à dérober.

Trois mesures répondent à ce risque.

\begin{itemize}[leftmargin=1.2cm]
  \item \textbf{Il n'est pas stocké en clair.} La base ne contient que son
  empreinte SHA-256. Une lecture de la base ne livre donc aucune session
  utilisable. Le hachage rapide suffit ici, contrairement au cas des mots de
  passe : le jeton est une valeur aléatoire de 48 octets, et non un secret
  choisi par un être humain --- il n'existe pas de dictionnaire de jetons.

  \item \textbf{Il tourne.} Chaque rafraîchissement révoque le jeton présenté
  et en émet un nouveau, rattaché à la même \emph{famille}.

  \item \textbf{Sa réutilisation est détectée.} Présenter deux fois le même
  jeton est impossible en usage normal. C'est en revanche exactement ce qui se
  produit lorsqu'un jeton a été dérobé : l'attaquant et la cliente légitime
  l'utilisent chacun de leur côté. Impossible de savoir lequel des deux se
  présente --- le système révoque donc \emph{toute la famille}. Les deux sont
  déconnectés, et l'accès frauduleux prend fin. Le coût pour la cliente est une
  reconnexion ; le bénéfice est la fin d'une session volée qui, sans cela, se
  prolongerait indéfiniment par rotations successives.
\end{itemize}

Côté mobile, le jeton n'est pas rangé dans le stockage ordinaire de
l'application mais dans le \textbf{coffre sécurisé du système
d'exploitation} --- trousseau iOS, \textit{Keystore} Android --- adossé au
matériel du téléphone.

Le déroulement complet, rotation et détection de réutilisation comprises, est
donné par la \figref{seq-refresh} du chapitre~2.

\subsection{Codes à usage unique}

La vérification d'adresse et la réinitialisation de mot de passe reposent sur
un code à six chiffres envoyé par courriel. Un tel code est court par
nécessité --- il doit être recopié à la main --- donc faible par construction :
un million de possibilités se parcourent en quelques minutes. Cinq mesures le
rendent néanmoins sûr.

\begin{table}[H]
\centering
\caption{Protection des codes à usage unique}
\label{tab:codes}
\begin{tabularx}{\textwidth}{L{4.0cm} X}
\toprule
\entete{Mesure} & \entete{Ce qu'elle empêche} \\
\midrule
Génération par \texttt{randomInt} &
Générateur cryptographique, et non \texttt{Math.random} dont la suite est
prédictible : un code observé ne permet pas de deviner le suivant. \\

Stockage haché &
Une lecture de la base ne livre aucun code exploitable. \\

Plafond de cinq tentatives &
Ferme la porte à l'énumération : sans lui, six chiffres se devinent. \\

Validité de quinze minutes &
Borne la fenêtre d'exploitation d'un code intercepté. \\

Comparaison en temps constant &
\texttt{timingSafeEqual} : le temps de réponse ne trahit pas le nombre de
chiffres corrects, ce qui interdit de reconstituer le code position par
position. \\
\bottomrule
\end{tabularx}
\end{table}

À quoi s'ajoute une propriété structurelle : la recherche s'effectue par
\emph{couple (compte, usage)} et jamais par le code seul. Un code deviné ne
vaut donc que pour le compte auquel il a été émis --- alors qu'une recherche
par code seul aurait fait qu'un code trouvé au hasard ouvre n'importe quel
compte l'ayant reçu, ce qui multiplie les chances de l'attaquant par le nombre
de codes en circulation. Enfin, émettre un nouveau code efface le précédent du
même usage : au plus un code actif par cliente et par usage.

\subsection{Défense contre l'énumération de comptes}

Savoir quelles adresses possèdent un compte est une information exploitable :
elle alimente le harponnage ciblé et, s'agissant d'un institut de beauté, elle
révèle une information personnelle en soi.

Deux canaux la trahissent, et tous deux ont été fermés.

\paragraph{Le canal des messages.} La demande de réinitialisation répond la
même chose que l'adresse existe ou non.

\paragraph{Le canal du temps.} Ce canal est plus subtil, et c'est précisément
ce qui le rend dangereux : il subsiste alors même que les messages sont
identiques. Vérifier un mot de passe avec Argon2 coûte plusieurs centaines de
millisecondes. Si le serveur répondait immédiatement pour une adresse inconnue
--- puisqu'il n'a alors aucun haché à comparer --- et après ce calcul pour une
adresse connue, la seule mesure du délai de réponse suffirait à trier les
adresses.

La connexion exécute donc, lorsque l'adresse est inconnue, une vérification
\emph{factice} contre le haché d'une valeur aléatoire. Elle échoue toujours :
seul le temps consommé compte.

\begin{lstlisting}[caption={Uniformisation du temps de réponse à la connexion},captionpos=b]
// Hash de reference, calcule une seule fois, sur une valeur aleatoire que
// personne ne connait. Il ne protege rien : il sert uniquement a faire
// travailler argon2 aussi longtemps qu'une vraie verification, pour qu'un
// email inconnu reponde au meme rythme.
private async burnPasswordComparison(candidate: string): Promise<void> {
  this.dummyHashPromise ??= argon2.hash(randomBytes(32).toString('hex'));
  const dummyHash = await this.dummyHashPromise;
  await argon2.verify(dummyHash, candidate).catch(() => false);
}
\end{lstlisting}

% ─────────────────────────────────────────────────────────────────
\section{Autorisation et contrôle d'accès}

\subsection{Modèle de contrôle d'accès}

Le contrôle d'accès repose sur un modèle par rôles (RBAC) à trois niveaux
strictement ordonnés. Le tableau ci-dessous donne la matrice des droits
effectivement appliquée par le serveur.

\begin{table}[H]
\centering
\caption{Matrice des droits par rôle}
\label{tab:matrice-droits}
\begin{tabularx}{\textwidth}{X C{1.6cm} C{1.7cm} C{2.1cm} C{1.8cm}}
\toprule
\entete{Opération} & \entete{Public} & \entete{Cliente} & \entete{Personnel} & \entete{Gérante} \\
\midrule
Consulter les catalogues actifs   & \checkmark & \checkmark & \checkmark & \checkmark \\
Réserver, commander pour soi      &            & \checkmark & \checkmark & \checkmark \\
Consulter \emph{ses} rendez-vous et commandes & & \checkmark & \checkmark & \checkmark \\
Consulter ceux de \emph{toutes}   &            &            & \checkmark & \checkmark \\
Réserver pour le compte d'une cliente &        &            & \checkmark & \checkmark \\
Changer un statut, honorer un bon &            &            & \checkmark & \checkmark \\
Créditer des points à la main     &            &            & \checkmark & \checkmark \\
Gérer le catalogue et les stocks  &            &            &            & \checkmark \\
Régler horaires, capacité, fermetures &        &            &            & \checkmark \\
Définir récompenses et paliers    &            &            &            & \checkmark \\
Gérer les comptes et les rôles    &            &            &            & \checkmark \\
Consulter le journal d'audit      &            &            &            & \checkmark \\
Exporter les données (Excel)      &            &            &            & \checkmark \\
Téléverser une image              &            &            &            & \checkmark \\
\bottomrule
\end{tabularx}
\end{table}

\subsection{Application par gardes successifs}

Cette matrice n'est pas une convention documentaire : elle est appliquée par
deux gardes que toute requête traverse, au sein d'une chaîne de filtres qu'elle
parcourt \emph{avant} d'atteindre la moindre ligne de logique métier.

\figProjet[0.85]{securite.png}{Chaîne de traitement de sécurité d'une requête}{securite}

Le premier vérifie la signature et la validité du jeton. Il est actif par
défaut : une route n'échappe à l'authentification que si elle est
\emph{explicitement} déclarée publique. Ce sens du défaut est essentiel ---
avec la convention inverse, une route ajoutée un jour de fatigue serait
ouverte, et rien ne le signalerait.

Le second compare le rôle du compte aux rôles exigés par la route. Il est
indépendant de tout ce que l'interface affiche ou masque : une route
d'administration atteinte par un autre moyen --- l'outil de développement du
navigateur, un client HTTP en ligne de commande, la documentation OpenAPI ---
est refusée par le serveur lui-même.

\subsection{Contrôle de propriété des ressources}

Le rôle ne suffit pas. Deux clientes portent le même rôle, et rien dans ce rôle
n'empêche l'une de demander la commande de l'autre : c'est la faille de type
\textit{Insecure Direct Object Reference}, l'une des plus répandues et des plus
faciles à exploiter --- il suffit de changer un identifiant dans une adresse.

Chaque route qui expose une ressource nominative vérifie donc, en plus du rôle,
que la ressource appartient bien au demandeur. La vérification est faite
\emph{à la lecture en base}, en filtrant sur le propriétaire, et non après
coup : une commande qui n'appartient pas à la cliente est traitée comme
introuvable, ce qui évite au passage de confirmer son existence.

Les identifiants sont par ailleurs des UUID et non des entiers séquentiels. Ce
n'est pas une protection en soi --- la vérification de propriété reste la seule
qui compte --- mais cela ôte à un attaquant la possibilité d'énumérer les
ressources en incrémentant un nombre, et empêche de déduire le volume
d'activité de l'institut d'un identifiant observé.

\subsection{Invariants d'administration}

Deux invariants protègent le compte de direction, et ils sont appliqués par le
serveur :

\begin{itemize}[leftmargin=1.2cm]
  \item \textbf{une seule administratrice à la fois} --- promouvoir un compte
  exige de vérifier qu'aucun autre ne porte déjà ce rôle ;
  \item \textbf{jamais zéro administratrice} --- rétrograder la dernière est
  refusé. Sans ce garde-fou, une fausse manipulation rendrait la gestion du
  centre définitivement inaccessible, sans autre recours qu'une intervention
  directe en base.
\end{itemize}

% ─────────────────────────────────────────────────────────────────
\section{Sécurité des entrées et des sorties}

\subsection{Validation stricte des entrées}

Un tube de validation global s'applique à toutes les routes, avec trois
réglages qui comptent :

\begin{itemize}[leftmargin=1.2cm]
  \item \textbf{liste blanche} --- seules les propriétés déclarées dans le
  contrat de la route sont conservées ;
  \item \textbf{rejet des propriétés inconnues} --- une propriété non déclarée
  fait \emph{échouer} la requête au lieu d'être silencieusement ignorée. C'est
  ce réglage qui ferme la porte à l'affectation en masse : soumettre un champ
  \texttt{role} ou \texttt{pointsBalance} dans une mise à jour de profil est
  rejeté, et non toléré ;
  \item \textbf{conversion typée} --- les valeurs sont converties dans le type
  attendu avant d'atteindre la logique métier, qui ne reçoit donc jamais une
  chaîne là où elle attend un nombre ou une date.
\end{itemize}

Les règles de validation sont déclarées au plus près des données, dans les
objets de transfert, et sont partagées : la politique de mot de passe, le
format des codes à six chiffres et celui des numéros de téléphone marocains
sont des validateurs uniques, réutilisés partout où la donnée apparaît.

\subsection{Injection SQL}

Toutes les requêtes passent par l'ORM, qui les émet sous forme paramétrée :
les valeurs ne sont jamais concaténées dans le texte de la requête.

Le projet compte \emph{une} requête SQL écrite à la main --- la prise du verrou
de réservation. Elle est écrite avec un gabarit étiqueté (\textit{tagged
template}), lui aussi paramétré, et sa seule valeur variable est une constante
du code, jamais une entrée utilisateur.

\subsection{Contrôle des fichiers téléversés}

Le téléversement d'images est le seul chemin par lequel un utilisateur écrit un
fichier sur le serveur --- fichier ensuite republié en HTTP. C'est donc le
point le plus sensible du système, et il est traité comme tel.

\begin{itemize}[leftmargin=1.2cm]
  \item \textbf{Réservé à la gérante.} Aucune cliente ne peut téléverser quoi
  que ce soit.

  \item \textbf{Reconnaissance par signature binaire.} Le type déclaré par le
  navigateur et l'extension du nom de fichier viennent tous deux du client et
  se falsifient en une ligne. Seuls les premiers octets du fichier disent ce
  qu'il \emph{est} : trois signatures sont acceptées --- JPEG, PNG, WebP ---
  et c'est la signature reconnue, jamais l'extension fournie, qui détermine
  l'extension enregistrée.

  \item \textbf{Le SVG est refusé.} C'est un refus délibéré, et non un oubli :
  le SVG est du XML susceptible de contenir du JavaScript. Servi depuis le
  domaine de l'application, il s'exécuterait avec les droits de celle-ci ---
  une faille XSS stockée, ouverte par un simple téléversement d'image.

  \item \textbf{Nom de fichier tiré au sort.} Le nom fourni par le client est
  jeté ; le fichier est enregistré sous un identifiant aléatoire. Cela ferme la
  traversée de répertoire (un nom du type \texttt{../../etc/passwd}) et
  l'écrasement d'un fichier existant.

  \item \textbf{Taille plafonnée à 5 Mo}, et un seul fichier par requête.

  \item \textbf{Cinq minutes de recul.} Les images sont servies avec un cache
  d'un an, marqué immuable --- ce qui est sûr précisément parce que le nom est
  tiré au sort : remplacer une photo crée un nouveau nom, jamais une nouvelle
  version du même. Il n'y a donc rien à invalider.
\end{itemize}

\subsection{Fuite d'information par les messages d'erreur}

Un message d'erreur technique renseigne un attaquant : noms de tables, noms de
colonnes, noms de contraintes, parfois valeurs. Un filtre global intercepte
donc toute exception avant qu'elle n'atteigne le client et applique trois
traitements :

\begin{enumerate}[leftmargin=1.2cm]
  \item les exceptions métier levées volontairement passent intactes --- leur
  message est destiné à l'utilisatrice ;
  \item les erreurs de base de données identifiées sont traduites en langage
  métier ;
  \item \textbf{tout le reste devient un message neutre}, le détail et la pile
  d'appels ne vivant que dans les journaux du serveur.
\end{enumerate}

Les journaux eux-mêmes ont été traités comme un risque : une erreur serveur y
consigne l'\emph{identifiant} de l'utilisateur, jamais son adresse
électronique --- assez pour rejouer l'incident, pas assez pour reconstituer un
fichier d'adresses à partir des journaux.

\subsection{Injection d'en-tête de courriel}

Le nom de l'expéditeur est inséré entre guillemets dans un en-tête
\texttt{From}. Un retour à la ligne y couperait l'en-tête en deux et
permettrait d'en injecter d'autres --- un \texttt{Bcc} vers une adresse
tierce, par exemple. Les caractères de rupture sont donc retirés avant
insertion. La valeur vient certes de la configuration et non d'un client, mais
un fichier d'environnement recopié de travers suffirait à produire l'effet, et
l'échec serait silencieux.

% ─────────────────────────────────────────────────────────────────
\section{Sécurité du transport et des en-têtes}

\subsection{Chiffrement du transport}

L'API n'est publiée que sur l'interface locale du serveur. Ce choix est une
mesure de sécurité et non une limitation : il rend \emph{impossible} de la
joindre sans passer par le reverse proxy placé devant, à qui revient la
terminaison HTTPS. Publier le port directement reviendrait à servir en clair
les jetons d'authentification, sur des réseaux Wi-Fi publics compris.

Côté mobile, un garde-fou interdit de configurer une adresse d'API non
chiffrée en dehors du développement local.

\subsection{En-têtes de sécurité HTTP}

Helmet retire l'en-tête qui annonçait la technologie du serveur --- une
information sans usage légitime, qui oriente le choix des exploits --- et pose
une douzaine d'en-têtes que le navigateur applique.

\begin{table}[H]
\centering
\caption{Principaux en-têtes de sécurité posés}
\label{tab:entetes}
\begin{tabularx}{\textwidth}{L{6.3cm} X}
\toprule
\entete{En-tête} & \entete{Rôle} \\
\midrule
\texttt{Strict-Transport-Security} & Impose HTTPS pour les visites suivantes,
et ferme la fenêtre de rétrogradation vers HTTP. \\
\texttt{Content-Security-Policy} & Restreint les sources de scripts : première
barrière contre le XSS. \\
\texttt{X-Content-Type-Options} & Interdit au navigateur de deviner le type
d'une ressource, donc d'exécuter comme script un fichier qui n'en est pas un. \\
\texttt{X-Frame-Options} & Interdit l'inclusion dans un cadre : protection
contre le détournement de clic. \\
\texttt{Cross-Origin-Resource-Policy} & Contrôle quelles origines peuvent
charger une ressource. \\
\bottomrule
\end{tabularx}
\end{table}

\begin{attention}[Un arbitrage assumé]
La politique de sécurité du contenu interdit par défaut les scripts en ligne.
Or l'interface de documentation OpenAPI s'initialise justement par un script en
ligne : l'activer telle quelle afficherait une page blanche. Elle est donc
désactivée \emph{uniquement là où cette documentation tourne}, c'est-à-dire en
développement. En production, la documentation est coupée et la politique
s'applique pleinement. L'arbitrage porte ainsi sur un environnement qui n'est
pas exposé, et non sur le serveur en ligne.
\end{attention}

\subsection{Contrôle des origines}

La politique CORS n'autorise que les adresses déclarées du back-office. Un site
tiers ne peut donc pas faire émettre à un navigateur des requêtes authentifiées
vers l'API.

\begin{attention}[Un piège de déploiement]
Cette liste doit contenir l'adresse réelle du back-office déployé. Non
renseignée, elle retombe sur l'adresse de développement : le back-office ne
peut plus joindre l'API, et l'échec ne se manifeste que par un blocage dans la
console du navigateur --- \emph{jamais} par un message du serveur. On cherche
donc longtemps du mauvais côté.
\end{attention}

% ─────────────────────────────────────────────────────────────────
\section{Disponibilité et résistance à l'abus}

\subsection{Limitation du débit}

Un compteur par adresse IP plafonne l'ensemble de l'API à 120 requêtes par
minute. Les routes coûteuses --- celles qui déclenchent un hachage Argon2 ou un
envoi de courriel --- portent en plus une limite stricte.

\begin{table}[H]
\centering
\caption{Limites de débit par catégorie de route}
\label{tab:rate-limit}
\begin{tabularx}{\textwidth}{L{5.4cm} C{3.0cm} X}
\toprule
\entete{Route} & \entete{Limite} & \entete{Ce qu'elle protège} \\
\midrule
Ensemble de l'API & 120 / min &
Saturation générale du service \\
Connexion, inscription & 10 / min &
Force brute sur les mots de passe ; épuisement du processeur par Argon2 \\
Changement de mot de passe & 10 / min & Force brute sur le mot de passe actuel \\
Vérification d'adresse & 10 / min & Énumération du code à six chiffres \\
Renvoi de code & 5 / min & Usage de l'API comme relais d'envoi de courriels \\
Mot de passe oublié & 5 / 5 min & Harcèlement par courriel d'une adresse tierce \\
Téléversement d'image & 30 / min & Saturation du disque du serveur \\
\bottomrule
\end{tabularx}
\end{table}

\begin{attention}[Le réglage qui décide de tout]
Le compteur s'appuie sur l'adresse IP du demandeur. Derrière un reverse proxy,
cette adresse est celle \emph{du proxy}, identique pour tout le monde : sans le
réglage correspondant, l'institut entier partage un unique compteur de 120
requêtes par minute, et l'API se met à répondre 429 au hasard dès qu'il y a du
monde. Le symptôme est déroutant --- parfait en test avec une seule personne,
cassé en production --- et rien dans les journaux ne le rattache à sa cause.

Le réglage inverse est tout aussi dangereux, et c'est ce qui en fait un vrai
arbitrage : activé sans proxy devant, l'en-tête d'adresse réelle devient
falsifiable par n'importe qui, et il suffit alors d'en changer la valeur à
chaque requête pour contourner entièrement la limitation. Ce réglage n'est donc
à activer \emph{que} si un proxy de confiance est réellement en place.
\end{attention}

\subsection{Plafonnement des volumes}

Deux mesures évitent qu'une requête légitime ne dégénère en déni de service.

\paragraph{La pagination obligatoire.} Toute liste susceptible de croître ---
historique de fidélité, commandes, rendez-vous, utilisateurs --- est paginée
côté serveur. Sans cela, une cliente fidèle depuis trois ans rechargeait la
totalité de son historique à chaque ouverture de l'écran.

\paragraph{Le plafond d'export.} Les exports Excel comptent les lignes
\emph{avant} de charger quoi que ce soit et refusent au-delà de 20 000 lignes,
en invitant à restreindre la période. Un export non borné aurait constitué
l'opération la plus lourde de l'API, accessible d'une seule requête. Le
comptage préalable importe autant que le plafond : c'est lui qui garantit que
le refus ne coûte rien.

\subsection{Correction sous concurrence}

Les mécanismes décrits au chapitre~2 --- verrou consultatif sur les
réservations, écritures conditionnelles atomiques sur le stock et sur le solde
de points --- relèvent aussi de la sécurité, et pas seulement de la
correction fonctionnelle.

Une condition de course est en effet exploitable dès lors qu'elle est
déclenchable à volonté. Le débit conditionnel du solde de points en est
l'illustration : lire le solde puis le décrémenter permettrait, en lançant
plusieurs échanges simultanés, de dépenser plusieurs fois les mêmes points et
d'obtenir des récompenses non acquises. Le débit est donc exprimé comme une
\emph{écriture conditionnelle} --- retirer $n$ points à condition qu'il y en ait
au moins $n$ --- évaluée par la base au moment de l'écriture. Si la condition
n'est pas satisfaite, aucune ligne n'est modifiée et rien n'est émis.

% ─────────────────────────────────────────────────────────────────
\section{Traçabilité et non-répudiation}

Certaines opérations ont une valeur économique et doivent laisser une trace
opposable.

\begin{itemize}[leftmargin=1.2cm]
  \item \textbf{L'ajout manuel de points} exige un motif et enregistre
  l'employée qui l'a saisi. Un journal d'audit, consultable par la seule
  gérante, liste ces écritures.

  \item \textbf{La remise d'un bon} enregistre la date et l'identité de la
  personne qui l'a honoré.

  \item \textbf{Le lien vers l'employée n'est jamais supprimé en cascade.} Le
  départ d'une employée ne doit ni effacer les points des clientes qu'elle a
  servies, ni la trace de ce qui leur a été remis : l'écriture demeure, elle ne
  désigne simplement plus personne.

  \item \textbf{Les montants sont figés} au moment de l'acte. Une révision de
  tarif ne réécrit pas rétroactivement ce qui a été facturé, ni les points
  gagnés.
\end{itemize}

% ─────────────────────────────────────────────────────────────────
\section{Protection des données personnelles}

\subsection{Cadre applicable}

Le traitement de données personnelles au Maroc est régi par la \textbf{loi
09-08}, dont la Commission Nationale de contrôle de la protection des Données à
caractère Personnel (CNDP) assure le contrôle. Trois obligations ont eu un
effet direct sur la conception.

\paragraph{La minimisation.} Seules sont collectées les données nécessaires :
identité, adresse électronique, téléphone --- ce dernier parce qu'un institut
doit pouvoir rappeler une cliente pour confirmer une commande. Aucune donnée de
paiement n'est traitée : le règlement s'effectue à la remise, en espèces, et
l'application n'en voit rien. C'est une réduction considérable du risque, et
elle découle d'un choix fonctionnel plutôt que technique.

\paragraph{Le droit d'accès.} La cliente consulte à tout moment son profil, ses
rendez-vous, ses commandes et l'intégralité de son historique de fidélité.

\paragraph{Le droit à l'effacement.} C'est celui qui a demandé le plus de
travail, et il fait l'objet de la section suivante.

\subsection{L'effacement par anonymisation}

Le droit à l'effacement se heurte ici à une exigence contradictoire : les
commandes et les rendez-vous d'une cliente portent le chiffre d'affaires de
l'institut, qui doit lui survivre. Une suppression en cascade détruirait la
comptabilité ; un refus de supprimer méconnaîtrait la loi.

La solution retenue est l'\textbf{anonymisation} : la ligne survit, vidée de
tout ce qui identifie la personne.

\begin{lstlisting}[caption={Anonymisation d'un compte : ce qui est effacé, ce qui survit},captionpos=b]
await this.prisma.$transaction([
  this.prisma.user.update({
    where: { id: userId },
    data: {
      // L'adresse d'origine est LIBEREE : si la cliente revient un jour, elle
      // pourra se reinscrire avec le meme email — sur un compte neuf. Le
      // domaine `.invalid` est reserve par la RFC 2606 et n'est routable nulle
      // part : aucun email ne partira jamais vers cette coquille.
      email: `supprime-${userId}@compte-supprime.invalid`,
      firstName: null, lastName: null, phone: null,
      passwordHash,          // hache d'un secret aleatoire que nul ne connait
      emailVerifiedAt: null,
      deletedAt: new Date(),
    },
  }),
  // Les sessions ouvertes tombent tout de suite.
  this.prisma.refreshToken.deleteMany({ where: { userId } }),
  // Et un code de reinitialisation en cours ne doit pas servir de porte
  // derobee sur une coquille anonyme.
  this.prisma.verificationCode.deleteMany({ where: { userId } }),
]);
\end{lstlisting}

Plusieurs points de ce traitement méritent d'être relevés.

\begin{itemize}[leftmargin=1.2cm]
  \item \textbf{L'opération est transactionnelle.} L'anonymisation, la
  destruction des sessions et celle des codes en cours réussissent ou échouent
  ensemble. Un échec partiel laisserait une coquille anonyme accessible par une
  session encore ouverte.

  \item \textbf{Le domaine \texttt{.invalid} est réservé par la RFC 2606} et
  n'est routable nulle part. Aucun courriel ne peut partir vers cette coquille,
  même par erreur de code --- ce qu'un domaine inventé ne garantirait pas.

  \item \textbf{L'adresse d'origine est libérée}, ce qui permet une
  réinscription ultérieure sur un compte neuf.

  \item \textbf{Trois verrous indépendants} ferment le compte : le refus de
  connexion sur le marqueur de suppression, le mot de passe devenu inconnu de
  tous, et le rejet à chaque requête des jetons d'accès déjà émis.

  \item \textbf{La garantie est structurelle.} Les rendez-vous et les commandes
  sont rattachés à la cliente en mode \texttt{Restrict} : si du code tentait un
  jour une suppression, la contrainte le ferait échouer bruyamment au lieu de
  détruire l'historique en silence.
\end{itemize}

L'opération est enfin \textbf{idempotente} : rejouée sur un compte déjà
anonymisé, elle répond succès sans rien réécrire --- ce qui préserve la date
d'origine, seule trace de la date à laquelle le droit a été exercé.

\subsection{Gestion des secrets}

Aucun secret ne figure dans le dépôt. Les fichiers d'environnement ne sont pas
versionnés ; seuls leurs modèles le sont, documentés valeur par valeur.

La configuration est \textbf{validée au démarrage}, et l'API refuse de démarrer
si elle est incohérente. Ce choix mérite d'être souligné : la solution
alternative --- démarrer et se débrouiller --- produit des états dégradés que
rien ne signale.

\begin{itemize}[leftmargin=1.2cm]
  \item \textbf{Le secret de signature doit compter au moins 32 caractères.} Un
  secret plus court se casse hors ligne, et permet ensuite de forger un jeton
  d'administration sans jamais toucher à l'API.

  \item \textbf{L'envoi réel de courriels est exigé en production.} Un serveur
  de production laissé en mode console n'enverrait aucun code de vérification
  ni de réinitialisation, et rien ne le signalerait : les clientes se
  retrouveraient simplement incapables de valider leur compte.

  \item \textbf{Les valeurs de repli sont bannies} là où elles seraient
  dangereuses. Le secret de signature est lu par une méthode qui lève une
  exception s'il manque, à la signature comme à la vérification.
\end{itemize}

\subsection{Sauvegardes}

Un script produit un cliché horodaté de la base et supprime ceux de plus de
quatorze jours. Deux règles d'exploitation l'accompagnent, et elles comptent
autant que le script : le cliché doit être recopié \emph{hors de la machine}
--- une sauvegarde qui vit sur le serveur qu'elle sauvegarde ne protège de rien
--- et le dossier des images, qui vit hors de la base, doit être copié
séparément. La procédure de restauration et sa vérification sont documentées
dans le script lui-même.

% ─────────────────────────────────────────────────────────────────
\section{Sécurité de la chaîne de développement}

La sécurité d'une application ne s'arrête pas à son code : elle englobe ses
dépendances, sa construction et son déploiement.

\paragraph{Les dépendances.} Les versions sont figées par un fichier de
verrouillage, et deux dépendances \emph{transitives} sont forcées à des
versions corrigées : une bibliothèque d'analyse YAML tirée par la couche de
documentation, et une bibliothèque d'identifiants tirée par le générateur de
classeurs Excel. Ces dépendances-là sont les plus faciles à négliger --- on ne
les a pas choisies, et elles n'apparaissent pas dans la liste des dépendances
directes.

\paragraph{L'image de production.} Elle est construite en trois étapes, si bien
que l'image livrée ne contient ni compilateur, ni cadriciel de test, ni
outillage de développement --- une réduction de surface d'attaque autant qu'un
gain de taille. Le processus y tourne sous un utilisateur non privilégié, et
non en tant que \texttt{root}.

\paragraph{L'intégration continue.} Chaque modification rejoue l'intégralité des
tests contre une véritable base PostgreSQL. Plusieurs de ces tests portent
directement sur des propriétés de sécurité : rotation et détection de
réutilisation des jetons, anonymisation des comptes, capacité des créneaux sous
concurrence, restriction des exports.

\paragraph{L'hygiène du dépôt.} Les notes de travail internes, qui décrivaient
des faiblesses avant leur correction, ont été retirées du suivi. Le retrait est
documenté avec sa limite --- il ne réécrit pas l'historique, et le contenu
demeure accessible dans les commits passés. Cette honnêteté a une valeur
opérationnelle : elle évite de tenir pour effacé ce qui ne l'est pas.

% ─────────────────────────────────────────────────────────────────
\section{Synthèse}

Le tableau ci-dessous récapitule, pour chaque menace du modèle, la
contre-mesure retenue et l'endroit où elle est appliquée.

\begin{table}[H]
\centering
\caption{Synthèse menaces $\rightarrow$ contre-mesures}
\label{tab:synthese-securite}
\begin{tabularx}{\textwidth}{L{4.2cm} X}
\toprule
\entete{Menace} & \entete{Contre-mesure} \\
\midrule
Force brute sur les mots de passe &
Argon2id ; 10 tentatives par minute et par IP ; politique de complexité
identique sur les trois chemins de saisie \\

Vol d'un jeton de session &
Jeton d'accès de 15 minutes ; rafraîchissement haché, tourné, avec révocation
de toute la famille en cas de réutilisation ; coffre sécurisé du téléphone \\

Devinette d'un code à 6 chiffres &
Générateur cryptographique ; stockage haché ; 5 tentatives ; 15 minutes de
validité ; comparaison en temps constant ; recherche par (compte, usage) \\

Énumération de comptes &
Réponses identiques et \emph{temps de réponse} identiques, par vérification
factice sur adresse inconnue \\

Élévation de privilège &
Garde de rôle côté serveur sur chaque route ; rôle relu en base à chaque
requête ; rejet des propriétés non déclarées \\

Lecture des données d'autrui &
Contrôle de propriété au niveau de la requête en base ; identifiants UUID \\

Jeton émis avant suppression &
Compte rechargé et marqueur de suppression vérifié à chaque requête \\

Injection SQL &
Requêtes paramétrées par l'ORM ; unique requête manuelle sans entrée
utilisateur \\

XSS par téléversement &
Reconnaissance par signature binaire ; SVG refusé ; nom aléatoire ; en-têtes
\texttt{nosniff} et CSP \\

Fuite du schéma de base &
Filtre d'exceptions global ; documentation OpenAPI coupée en production \\

Déni de service &
Limitation de débit globale et par route ; pagination obligatoire ; plafond
d'export à 20 000 lignes \\

Condition de course exploitable &
Verrou consultatif sur les réservations ; écritures conditionnelles atomiques
sur le stock et le solde \\

Interception du trafic &
API sur interface locale ; HTTPS imposé par le reverse proxy ; HSTS ;
garde-fou HTTPS côté mobile \\

Répudiation d'une opération &
Motif obligatoire et auteur enregistré ; journal d'audit ; liens vers
l'employée jamais supprimés en cascade \\

Atteinte aux données personnelles &
Minimisation ; aucune donnée de paiement ; anonymisation transactionnelle ;
sauvegardes horodatées \\
\bottomrule
\end{tabularx}
\end{table}

% ─────────────────────────────────────────────────────────────────
\section{Limites connues et durcissement à venir}

Aucun système n'est sûr dans l'absolu, et un rapport qui prétendrait le
contraire manquerait son objet. Les limites suivantes sont assumées et
identifiées.

\paragraph{L'absence de second facteur.} L'authentification repose sur un seul
facteur. Pour un compte cliente, l'enjeu reste mesuré ; pour le compte de
direction, qui donne accès à l'ensemble du fichier clientes, un second facteur
serait la première amélioration à apporter.

\paragraph{Les données au repos ne sont pas chiffrées.} La base repose sur le
chiffrement du disque du serveur hôte. Un chiffrement au niveau des colonnes
les plus sensibles constituerait une couche supplémentaire.

\paragraph{Le compteur de débit vit en mémoire.} Il est donc propre à chaque
instance et remis à zéro au redémarrage. Tant que l'API tourne en un seul
exemplaire, la protection est effective ; une mise à l'échelle horizontale
imposerait de déporter ce compteur dans un magasin partagé.

\paragraph{Aucun audit externe n'a été mené.} Les propriétés de sécurité sont
établies par les tests automatisés et par la revue du code. Un test
d'intrusion mené par un tiers reste la seule façon de découvrir ce que la
personne qui a écrit le code ne pense pas à chercher.

\paragraph{La supervision est minimale.} Les événements sont journalisés, mais
rien ne les agrège ni ne déclenche d'alerte. Une série d'échecs de connexion
sur un même compte devrait alerter, et ne le fait pas aujourd'hui.

\paragraph{Les points restant à faire avant l'ouverture publique.} Le
certificat HTTPS et le reverse proxy conditionnent plusieurs des mesures
décrites ici : tant qu'ils ne sont pas en place, la protection du transport
repose sur le fait que le service n'est pas encore exposé.

\section*{Conclusion}
\addcontentsline{toc}{section}{Conclusion}

Ce chapitre a présenté la sécurité de la plate-forme comme une chaîne
continue, du modèle de menaces jusqu'aux limites qui subsistent. Trois idées
en résument la démarche.

D'abord, \textbf{la frontière de confiance passe à l'entrée de l'API}. Ni
l'application mobile, ni le back-office, ni aucune valeur venue d'un client
n'est cru sur parole ; ce qui est masqué dans une interface n'est jamais
protégé pour autant.

Ensuite, \textbf{les défauts sont fermés}. Une route est authentifiée et
comptée sauf déclaration contraire, une propriété inconnue fait échouer la
requête, et une configuration incohérente empêche le démarrage plutôt que de
produire un état dégradé silencieux.

Enfin, \textbf{la sécurité n'a pas été ajoutée après coup}. Les mécanismes qui
garantissent qu'un créneau n'est pas vendu deux fois sont les mêmes qui
empêchent de dépenser deux fois les mêmes points ; le figement des montants
sert à la fois l'exactitude comptable et la non-répudiation ; l'anonymisation
concilie un droit légal et la survie de l'historique. Correction fonctionnelle
et sécurité se sont révélées être, le plus souvent, deux vues du même
problème.

Le chapitre suivant présente la mise en \oe uvre effective de la solution et
les tests qui en établissent les propriétés.

% ═══════════════════════════════════════════════════════════════════
%  Chapitre 4 — Mise en œuvre
% ═══════════════════════════════════════════════════════════════════
\chapter{Mise en \oe uvre}

\section*{Introduction}
\addcontentsline{toc}{section}{Introduction}

Ce chapitre présente la réalisation effective de la solution. Il décrit d'abord
la structure du code des trois applications, puis la stratégie de tests mise en
place et son automatisation par intégration continue. Il présente ensuite les
interfaces développées, illustrées par des captures d'écran, et se termine par
les modalités de déploiement et d'exploitation de la plate-forme.

% ─────────────────────────────────────────────────────────────────
\section{Structure du projet}

Le projet est organisé en mono-dépôt : les trois applications vivent dans un
même dépôt Git, à trois emplacements distincts. Ce choix évite la
désynchronisation entre le contrat de l'API et ses consommateurs --- une
modification du serveur et son adaptation côté client sont validées ensemble
par la même chaîne d'intégration continue.

\begin{table}[H]
\centering
\caption{Organisation du mono-dépôt}
\label{tab:monodepot}
\begin{tabularx}{\textwidth}{L{2.4cm} L{3.4cm} X}
\toprule
\entete{Dossier} & \entete{Rôle} & \entete{Pile technique} \\
\midrule
\texttt{backend/} & API REST, règles métier, base de données &
NestJS · Prisma · PostgreSQL \\
\texttt{backoffice/} & Interface de gestion &
React · Vite · Zustand \\
\texttt{mobile/} & Application cliente &
Expo · React Native · i18next \\
\texttt{scripts/} & Sauvegarde de la base & Shell \\
\texttt{brand/} & Identité visuelle et génération des icônes & Python \\
\bottomrule
\end{tabularx}
\end{table}

\subsection{Le backend}

Le code de l'API représente environ 8~000 lignes, auxquelles s'ajoutent 5~150
lignes de tests. Il est découpé en quatorze modules NestJS, réunis par un
module racine.

\figProjet[0.7]{structure-backend.png}{Structure des modules du backend}{structure-backend}

Douze d'entre eux exposent un contrôleur et suivent tous la même organisation,
ce qui rend le code prévisible :

\begin{itemize}[leftmargin=1.2cm]
  \item le \textbf{contrôleur} déclare les routes, leurs droits d'accès, leurs
  limites de débit et leur documentation ; il ne contient aucune règle ;
  \item le \textbf{service} porte la logique métier et les accès aux données ;
  \item les \textbf{objets de transfert} décrivent et valident les entrées ;
  \item les \textbf{tests} accompagnent le module qu'ils vérifient.
\end{itemize}

Les deux modules restants n'exposent aucune route : \texttt{prisma} met le
client de base de données à la disposition des autres, et \texttt{mail}
l'acheminement des courriels.

Trois dossiers, enfin, ne sont pas des modules mais du code partagé :
\texttt{common} regroupe le filtre d'exceptions, l'intercepteur de traduction,
la pagination et les validateurs communs, \texttt{config} valide les variables
d'environnement au démarrage, et \texttt{i18n} porte les messages traduits de
l'API.

\paragraph{Le c\oe ur du système.} Trois services concentrent l'essentiel de la
complexité : le service des rendez-vous (693 lignes), qui porte le moteur de
créneaux, la gestion du fuseau et le verrou de réservation ; le service de
fidélité (837 lignes), qui porte les deux circuits de récompense et
l'idempotence des crédits ; et le service d'authentification (477 lignes), qui
porte la rotation des jetons et les défenses contre l'énumération de comptes.

\subsection{Le back-office}

Le back-office représente environ 6~500 lignes de TypeScript et React,
organisées en huit pages et sept composants partagés.

\figProjet[0.7]{structure-backoffice.png}{Structure du back-office}{structure-backoffice}

\begin{itemize}[leftmargin=1.2cm]
  \item \texttt{pages/} --- une page par domaine métier : tableau de bord,
  commandes, rendez-vous, catalogue, fidélité, horaires, utilisateurs,
  connexion ;
  \item \texttt{components/} --- les éléments réutilisés : mise en page,
  pagination, boîte de confirmation, garde de route authentifiée, formulaire de
  catalogue, fenêtre d'export, fenêtre de réservation pour une cliente ;
  \item \texttt{api/} --- une façade par domaine, seule à connaître les
  adresses des routes ; les pages ne manipulent jamais d'URL ;
  \item \texttt{stores/} --- l'état global, réduit à la session authentifiée ;
  \item \texttt{theme/} --- la charte graphique partagée.
\end{itemize}

Le client HTTP centralise deux comportements délicats : l'ajout automatique du
jeton d'accès à chaque requête, et le rafraîchissement transparent de la
session lorsqu'il expire --- la requête ayant échoué est alors rejouée, sans
que l'utilisatrice s'aperçoive de quoi que ce soit.

\subsection{L'application mobile}

L'application mobile représente environ 9~700 lignes, réparties en vingt-quatre
écrans et dix-huit composants.

\figProjet[0.7]{structure-mobile.png}{Structure de l'application mobile}{structure-mobile}

La navigation s'organise en une barre d'onglets --- Accueil, Prestations,
Produits, Fidélité, Profil --- surmontant des piles de navigation propres à
chaque section. Le panier et la session sont conservés dans des magasins
d'état, le jeton de session étant stocké dans le coffre sécurisé du système
d'exploitation plutôt qu'en clair.

L'application est traduite en trois langues, avec chargement de la langue du
téléphone au démarrage. La prise en charge de l'arabe a demandé une attention
particulière : au-delà de la traduction, elle impose de vérifier chaque écran
dans le sens d'écriture de droite à gauche.

% ─────────────────────────────────────────────────────────────────
\section{Tests et validation}

\subsection{Stratégie de tests}

La stratégie retenue repose sur un constat simple : \emph{les propriétés les
plus importantes de ce projet ne sont pas démontrables par un test unitaire
classique}. Qu'un calcul de points donne le bon résultat se vérifie avec une
base simulée. Que trois réservations simultanées ne puissent pas occuper deux
cabines ne se vérifie qu'avec une vraie base, un vrai verrou et de vraies
transactions.

Les tests se répartissent donc en deux familles.

\begin{itemize}[leftmargin=1.2cm]
  \item \textbf{Les tests unitaires}, sur base simulée. Ils vérifient les
  règles métier pures : calcul des points, transitions d'états autorisées,
  génération de la grille de créneaux, validation des entrées, contrôle des
  droits.

  \item \textbf{Les tests d'intégration}, sur une véritable base PostgreSQL.
  Ils vérifient ce qu'une base simulée ne saurait démontrer : sérialisation des
  réservations concurrentes, atomicité du crédit de fidélité, annulation
  effective d'une transaction en échec, idempotence de l'émission des bons,
  rotation et détection de réutilisation des jetons, anonymisation des comptes.
\end{itemize}

Au total, le backend est couvert par 231 tests répartis en 25 suites, dont 8
écrivent réellement en base.

\begin{table}[H]
\centering
\caption{Répartition des suites de tests par domaine}
\label{tab:suites-tests}
\begin{tabularx}{\textwidth}{L{3.0cm} C{1.3cm} C{1.9cm} X}
\toprule
\entete{Domaine} & \entete{Suites} & \entete{Base réelle} & \entete{Propriétés vérifiées} \\
\midrule
Rendez-vous & 5 & 1 &
Grille de créneaux, fuseau du centre, fermetures, machine à états, capacité
sous concurrence \\

Fidélité & 4 & 2 &
Calcul et idempotence des crédits, débit conditionnel du solde, paliers de
visites, émission et remise des bons \\

Commandes & 5 & 2 &
Contrôle et décrément atomique du stock, restitution à l'annulation, crédit à
la remise, machine à états \\

Authentification & 3 & 1 &
Hachage, rotation des jetons, détection de réutilisation, validation du jeton
d'accès \\

Utilisateurs & 3 & 1 &
Profil, unicité de l'administratrice, anonymisation du compte supprimé \\

Exports & 2 & 0 &
Génération des classeurs, restriction à la gérante, bornes de dates \\

Horaires & 1 & 0 & Refus des plages qui se chevauchent \\
Infrastructure & 2 & 1 & Connexion à la base, sonde de santé \\
\midrule
\textbf{Total} & \textbf{25} & \textbf{8} & \textbf{231 tests} \\
\bottomrule
\end{tabularx}
\end{table}

\subsection{Le test de concurrence, en détail}

Le test le plus instructif du projet mérite d'être décrit, car il a
directement dicté une décision d'architecture.

Il place le centre à une capacité de deux places, puis lance trois réservations
\emph{réellement simultanées} sur le même créneau. La propriété attendue est
que deux réussissent et qu'une échoue --- jamais trois.

\begin{lstlisting}[caption={Principe du test de concurrence sur les réservations},captionpos=b]
// Trois réservations lancées en parallèle sur le même créneau,
// avec une capacité de 2 places.
const resultats = await Promise.allSettled([
  service.create(cliente1.id, { serviceId, startAt }),
  service.create(cliente2.id, { serviceId, startAt }),
  service.create(cliente3.id, { serviceId, startAt }),
]);

const reussies = resultats.filter((r) => r.status === 'fulfilled');
expect(reussies).toHaveLength(2);   // et jamais 3
\end{lstlisting}

Sur l'implémentation initiale --- compter les rendez-vous du créneau, puis
insérer --- ce test échouait de façon reproductible : les trois réservations
étaient acceptées. Les trois comptaient avant qu'aucune n'ait inséré, et les
trois trouvaient de la place. C'est ce test, et non une revue de code, qui a
imposé le recours au verrou consultatif décrit au chapitre~2.

Ces tests partagent la même base et le même verrou : ils doivent donc être
exécutés en série, ce qui est imposé par une option dédiée du lanceur de tests.

\subsection{Intégration continue}

L'ensemble est rejoué automatiquement à chaque modification par une chaîne
GitHub Actions. Elle comporte trois travaux.

\begin{table}[H]
\centering
\caption{Travaux de la chaîne d'intégration continue}
\label{tab:ci}
\begin{tabularx}{\textwidth}{L{2.8cm} X}
\toprule
\entete{Travail} & \entete{Étapes} \\
\midrule
Backend &
Installation des dépendances, génération du client Prisma, application des
migrations sur une base PostgreSQL éphémère, exécution du script d'amorçage,
compilation, puis exécution des 25 suites de tests. \\

Back-office & Installation des dépendances et vérification de types. \\

Mobile & Installation des dépendances et vérification de types. \\
\bottomrule
\end{tabularx}
\end{table}

Un point de cette chaîne mérite d'être signalé, car il n'est pas évident : le
script d'amorçage y est indispensable. Il crée les horaires d'ouverture, et
sans une seule ligne dans cette table, le centre est considéré fermé tous les
jours --- toute la suite de tests des rendez-vous échouerait pour une raison
qui n'a rien à voir avec le code testé.

\figProjet[0.85]{ci-github.png}{Exécution de la chaîne d'intégration continue}{ci-github}

\subsection{Résultats}

\figProjet[0.85]{resultats-tests.png}{Résultat de l'exécution complète de la suite de tests}{resultats-tests}

L'exécution complète est verte : 25 suites, 231 tests, aucun échec. Le
back-office et l'application mobile, qui n'ont pas de tests automatisés, sont
couverts par la vérification de types --- un filet plus mince, mais qui détecte
effectivement toute divergence entre le contrat de l'API et ses consommateurs.

% ─────────────────────────────────────────────────────────────────
\section{Interfaces graphiques}

\subsection{Application mobile}

\subsubsection{Authentification}

\figDouble{mobile-connexion.png}{Écran de connexion}{mobile-inscription.png}{Écran d'inscription}{mobile-auth}

L'inscription recueille l'adresse, le mot de passe, le nom, le prénom et le
téléphone. Ces trois derniers champs sont obligatoires : un rendez-vous ou une
commande sans nom ni numéro de téléphone est inexploitable pour l'institut, qui
doit pouvoir rappeler la cliente. Un code de vérification est envoyé par
e-mail à l'issue de l'inscription.

\figDouble{mobile-verification.png}{Vérification de l'adresse}{mobile-mdp-oublie.png}{Mot de passe oublié}{mobile-verif}

\subsubsection{Accueil et catalogues}

\figDouble{mobile-accueil.png}{Écran d'accueil}{mobile-prestations.png}{Catalogue des prestations}{mobile-accueil-prestations}

L'accueil réunit en une vue le solde de fidélité, une sélection de prestations
et une sélection de produits. C'est l'écran le plus consulté, et celui dont le
coût en requêtes a été le plus surveillé : il ne demande que ce qu'il affiche.

\figDouble{mobile-detail-prestation.png}{Détail d'une prestation}{mobile-produits.png}{Catalogue des produits}{mobile-details}

\subsubsection{Réservation}

\figDouble{mobile-reservation.png}{Choix de la date et du créneau}{mobile-recap-reservation.png}{Récapitulatif avant validation}{mobile-reservation}

L'écran de réservation affiche la grille des créneaux du jour choisi. Les
créneaux indisponibles --- capacité atteinte, ou heure déjà passée --- restent
visibles mais désactivés, plutôt que masqués : une cliente qui ne voit pas le
créneau de 15\,h ignore s'il est pris ou si l'institut n'ouvre pas à cette
heure-là.

Un jour fermé, qu'il le soit par les horaires hebdomadaires ou par une
fermeture exceptionnelle, est annoncé explicitement.

\figDouble{mobile-confirmation-rdv.png}{Confirmation du rendez-vous}{mobile-mes-rdv.png}{Mes rendez-vous}{mobile-rdv}

\subsubsection{Commande de produits}

\figDouble{mobile-detail-produit.png}{Détail d'un produit}{mobile-panier.png}{Panier}{mobile-produit-panier}

\figDouble{mobile-checkout.png}{Validation de commande}{mobile-mes-commandes.png}{Mes commandes}{mobile-commande}

La validation propose le retrait à l'institut ou la livraison à domicile,
laquelle exige une adresse. Le paiement s'effectue à la remise. La cliente peut
annuler sa commande tant qu'elle ne lui a pas été remise.

\subsubsection{Fidélité}

\figDouble{mobile-fidelite.png}{Écran de fidélité}{mobile-recompenses.png}{Mes récompenses et bons}{mobile-fidelite}

L'écran de fidélité réunit le solde de points, la progression vers le prochain
palier de visites, le catalogue des récompenses échangeables et l'historique
des mouvements. L'écran des récompenses présente les bons dus : ceux offerts
par un palier et ceux obtenus par échange de points, avec leur code à présenter
au comptoir.

\subsubsection{Profil}

\figDouble{mobile-profil.png}{Profil et réglages}{mobile-edition-profil.png}{Modification du profil}{mobile-profil}

Les réglages permettent de changer de langue, de basculer entre les modes clair
et sombre, de modifier le mot de passe et de supprimer le compte.

\subsection{Back-office}

\figProjet{bo-connexion.png}{Back-office --- écran de connexion}{bo-connexion}

Une cliente est refusée à l'entrée du back-office : ce contrôle est effectué
par le serveur et non par l'interface.

\figProjet{bo-dashboard.png}{Back-office --- tableau de bord}{bo-dashboard}

Le tableau de bord réunit ce qui appelle une action immédiate : les commandes
en attente de confirmation, les rendez-vous du jour et les produits en stock
faible. La notion d'« aujourd'hui » y est calculée par le serveur, dans le
fuseau du centre : c'est lui qui fait autorité, et non le poste qui consulte.

\figProjet{bo-commandes.png}{Back-office --- gestion des commandes}{bo-commandes}

Les commandes sont présentées par statut. Chaque changement de statut est
soumis à la machine à états : les seules transitions proposées sont celles que
le serveur accepterait.

\figProjet{bo-rendezvous.png}{Back-office --- gestion des rendez-vous}{bo-rendezvous}

Cette page permet également de réserver pour le compte d'une cliente ---
c'est-à-dire de saisir un rendez-vous pris par téléphone --- et de reporter un
rendez-vous existant, avec les mêmes contrôles de disponibilité que ceux
appliqués à une réservation cliente.

\figProjet{bo-catalogue.png}{Back-office --- gestion du catalogue}{bo-catalogue}

Le catalogue couvre les prestations, les produits et leurs catégories, y
compris le téléversement des photos et le suivi des stocks. Un élément retiré
du catalogue est désactivé et non supprimé : il disparaît immédiatement des
applications clientes, mais reste rattaché à l'historique des commandes et des
rendez-vous qui le référencent.

\figProjet{bo-fidelite.png}{Back-office --- programme de fidélité}{bo-fidelite}

Cette page réunit le catalogue des récompenses, les paliers de visites, la
liste des bons dus au comptoir, l'ajout manuel de points --- motif obligatoire
--- et le journal d'audit de ces ajouts, consultable par la seule gérante.

\figProjet{bo-horaires.png}{Back-office --- horaires et fermetures}{bo-horaires}

Les horaires se règlent par plages, à raison de plusieurs plages par jour. Un
jour sans aucune plage est fermé. Les fermetures exceptionnelles --- congés,
jours fériés --- priment sur les horaires hebdomadaires. La capacité du centre
et l'écart entre créneaux se règlent depuis cette même page, sans
redéploiement.

\figProjet{bo-utilisateurs.png}{Back-office --- gestion des utilisateurs}{bo-utilisateurs}

\figProjet[0.7]{bo-export.png}{Back-office --- fenêtre d'export Excel}{bo-export}

Les exports produisent des classeurs Excel des clientes, des commandes et des
rendez-vous, sur une période choisie. Le fichier est destiné à la gérante et
non à un développeur : aucune valeur technique brute n'y apparaît --- les
statuts y sont écrits en toutes lettres --- et les montants y restent des
\emph{nombres} plutôt que du texte, de sorte qu'une somme ou un tableau croisé
fonctionne dans le classeur.

\subsection{Courriers électroniques}

\figProjet[0.7]{mail-verification.png}{Courriel de vérification d'adresse}{mail-verification}

Les courriels sont envoyés dans la langue choisie par la cliente. En
développement, aucun message ne part : les codes s'affichent dans les journaux
du serveur, ce qui permet de dérouler l'intégralité des parcours --- inscription,
vérification, mot de passe oublié --- sans dépendre d'un fournisseur de
messagerie.

% ─────────────────────────────────────────────────────────────────
\section{Déploiement et exploitation}

\subsection{Mise en production}

La pile complète se lance par une commande unique, qui enchaîne la base de
données, l'application des migrations et le démarrage de l'API. La
documentation interactive de l'API est automatiquement coupée en production.

\subsection{Sauvegardes}

Un script produit un cliché horodaté de la base et supprime ceux de plus de
quatorze jours, à automatiser par une tâche planifiée. Les deux règles
d'exploitation qui l'accompagnent --- recopier le cliché hors de la machine, et
copier séparément le dossier des images, qui vit hors de la base --- ont été
exposées au chapitre~4 : elles relèvent de la disponibilité des données autant
que de l'exploitation courante.

La procédure de restauration, et la vérification qui doit la suivre, sont
documentées dans le script lui-même --- à l'endroit où on la cherchera le jour
où elle servira.

\subsection{Ce qui reste à faire avant une ouverture publique}

Par honnêteté sur l'état de la livraison, trois points restent ouverts :

\begin{itemize}[leftmargin=1.2cm]
  \item l'acquisition d'un nom de domaine et d'un certificat HTTPS ;
  \item la mise en place d'un serveur de messagerie définitif, à la place de la
  solution provisoire actuelle ;
  \item les notifications \textit{push}, qui dépendent d'une compilation
  native, donc du nom de domaine.
\end{itemize}

\section*{Conclusion}
\addcontentsline{toc}{section}{Conclusion}

Ce chapitre a présenté la mise en \oe uvre de la solution : la structure du code
des trois applications, la stratégie de tests et son automatisation, les
interfaces réalisées et les modalités d'exploitation. La démonstration du
produit final à travers ses écrans montre une plate-forme complète et
cohérente ; les 231 tests qui l'accompagnent, dont ceux qui écrivent dans une
véritable base, en établissent les propriétés les plus difficiles.

% ═══════════════════════════════════════════════════════════════════
%  Conclusion générale et perspectives
% ═══════════════════════════════════════════════════════════════════
\pageLiminaire{Conclusion Générale et Perspectives}

\section*{Bilan du travail réalisé}

Ce projet de fin d'études a abouti à une plate-forme logicielle complète pour
\yani\ : une API REST portant l'intégralité des règles métier, une application
mobile trilingue destinée aux clientes, et un back-office web destiné à
l'institut. Les trois applications partagent une même base PostgreSQL et un
même langage, et sont validées ensemble par une chaîne d'intégration continue.

Les objectifs fixés au départ ont été atteints. La réservation est possible à
toute heure, sur une grille de créneaux calculée à partir des horaires réels du
centre. Le carnet papier est remplacé par une base sauvegardée, consultable
depuis plusieurs postes. La carte de fidélité tamponnée a laissé place à un
compte doté d'un solde, d'un historique, de paliers de visites et de bons
traçables. Le stock est suivi, contrôlé et décrémenté au bon moment. La gérante
dispose enfin d'une visibilité sur son activité et de l'export de ses données.

Le travail ne s'est cependant pas limité à la mise en \oe uvre de
fonctionnalités. Sa part la plus exigeante a porté sur des propriétés qui ne se
voient pas à l'écran :

\begin{itemize}[leftmargin=1.2cm]
  \item \textbf{la correction sous concurrence} --- un créneau n'est jamais
  vendu deux fois, une unité de stock jamais attribuée deux fois, un point de
  fidélité jamais crédité deux fois ; ces propriétés sont établies par des
  tests qui écrivent dans une véritable base PostgreSQL, parce qu'aucune base
  simulée ne saurait les démontrer ;

  \item \textbf{la fidélité au temps réel} --- les instants sont stockés en
  temps universel, les horaires exprimés en heure locale du centre, et la
  conversion n'est faite qu'à un seul endroit, par le serveur, qui possède la
  base des fuseaux ;

  \item \textbf{la conservation de l'historique} --- une cliente peut exercer
  son droit à l'effacement sans que le chiffre d'affaires qu'elle a généré ne
  disparaisse avec elle ; les modes de suppression déclarés dans le schéma
  rendent cette garantie structurelle plutôt que déclarative ;

  \item \textbf{la sécurité} --- traitée à partir d'un modèle de menaces
  explicite et non comme une couche ajoutée en fin de parcours : hachage
  Argon2id, jetons courts, rotation avec révocation de famille en cas de
  réutilisation détectée, uniformisation du temps de réponse contre
  l'énumération de comptes, autorisation vérifiée par le serveur sur chaque
  route et doublée d'un contrôle de propriété des ressources, validation
  stricte des entrées, reconnaissance des fichiers téléversés par leur
  signature binaire.
\end{itemize}

\section*{Apports personnels}

Sur le plan technique, ce projet a été l'occasion de mener seul, de bout en
bout, une réalisation à trois cibles : un serveur, une application web et une
application mobile. Il m'a permis d'approfondir des sujets que le cursus
aborde nécessairement de façon théorique --- la concurrence en base de données,
la gestion des fuseaux horaires, la sécurité applicative --- et de découvrir
qu'ils ne se comprennent réellement qu'au moment où un test les met en échec.

L'apport le plus net concerne la sécurité, et il tient moins à des techniques
qu'à un changement de regard. Appliquer une grille de menaces à un système que
l'on a soi-même écrit oblige à l'examiner du point de vue de quelqu'un qui
cherche à en abuser, et non de celui qui l'utilise comme prévu. C'est cet
exercice qui a fait apparaître ce qu'aucune relecture fonctionnelle ne
signalait : que masquer un bouton n'a jamais protégé la route correspondante,
qu'un temps de réponse suffit à trahir l'existence d'un compte, et qu'une
condition de course déclenchable à volonté n'est pas un défaut de correction
mais une vulnérabilité.

Le test de concurrence sur les réservations en est l'illustration la plus
nette. L'implémentation initiale --- compter puis insérer --- paraissait
évidemment correcte à la lecture. Il a fallu trois réservations réellement
simultanées pour montrer qu'elle ne l'était pas. C'est l'enseignement que je
retiens le plus fermement de ce projet : \emph{en matière de concurrence,
l'intuition ne vaut rien, et seul un test qui reproduit la simultanéité fait
autorité}.

Sur le plan méthodologique, la démarche agile a démontré son intérêt de façon
concrète. Plusieurs des règles les plus importantes du système --- qu'un
rendez-vous occupe un créneau fixe et non la durée de sa prestation, qu'une
réclamation de récompense doive laisser une trace au comptoir, qu'un produit
retiré du catalogue ne doive pas s'effacer de l'historique --- ne figuraient
dans aucune spécification initiale. Elles sont apparues en montrant des
versions fonctionnelles à la personne qui allait s'en servir.

Sur le plan personnel, travailler pour une utilisatrice unique, identifiée et
non informaticienne a été formateur. Un message d'erreur incompréhensible n'est
pas un détail lorsque la personne qui le reçoit n'a personne à qui le
transmettre. Cette contrainte a orienté un grand nombre de décisions, du
libellé des statuts dans les exports Excel jusqu'au choix d'afficher les
créneaux indisponibles plutôt que de les masquer.

\section*{Difficultés rencontrées}

Trois difficultés méritent d'être mentionnées, ainsi que la façon dont elles
ont été surmontées.

\paragraph{La concurrence sur les réservations.} Décrite plus haut, elle a
imposé de descendre au niveau du SGBD et de recourir à un verrou consultatif
PostgreSQL --- une primitive qui ne figure dans aucune abstraction d'ORM.

\paragraph{Les fuseaux horaires.} Les erreurs de fuseau sont silencieuses :
elles ne lèvent aucune exception et produisent des résultats plausibles mais
faux, décalés d'une heure. La convention a donc été fixée explicitement ---
stockage en temps universel, horaires en heure locale, conversion en un seul
endroit --- et documentée dans le code comme dans l'API.

\paragraph{Les pièges de déploiement invisibles.} Plusieurs réglages
d'exploitation ne cassent rien immédiatement, mais produisent des symptômes
déroutants une fois en service : le réglage du proxy de confiance, sans lequel
tout l'institut partage un seul compteur de requêtes ; la liste des origines
autorisées, dont l'oubli n'apparaît que dans la console du navigateur. Ces
pièges ont été documentés à l'endroit où ils se manifestent, avec leur symptôme
et leur cause, plutôt que dans une documentation séparée que personne ne
consulte au moment utile.

\section*{Perspectives}

La plate-forme livrée est fonctionnelle, mais plusieurs évolutions ont été
identifiées, classées ici par horizon.

\paragraph{À court terme --- finaliser la mise en service.}
\begin{itemize}[leftmargin=1.2cm,nosep]
  \item Acquisition d'un nom de domaine et d'un certificat HTTPS, avec le
  reverse proxy correspondant.
  \item Mise en place d'un serveur de messagerie définitif.
  \item Publication de l'application sur les magasins Android et iOS.
\end{itemize}

\paragraph{À moyen terme --- enrichir l'usage.}
\begin{itemize}[leftmargin=1.2cm,nosep]
  \item \textbf{Notifications \textit{push}} : rappel de rendez-vous la veille,
  avis de commande prête, alerte de palier de fidélité franchi. C'est
  l'évolution la plus demandée, et elle dépend d'une compilation native, donc
  du nom de domaine.
  \item \textbf{Paiement en ligne}, en complément du paiement à la remise.
  \item \textbf{Statistiques de pilotage} : chiffre d'affaires par période,
  prestations les plus demandées, taux de retour de la clientèle. Les données
  sont déjà en base ; il ne manque que leur restitution.
  \item \textbf{Gestion fine des ressources} : aujourd'hui, la capacité du
  centre est un nombre unique de places réservables. Distinguer les cabines et
  les esthéticiennes permettrait de proposer des créneaux différenciés selon la
  prestation.
\end{itemize}

\paragraph{À moyen terme --- durcir.} Les limites de sécurité inventoriées au
chapitre~4 appellent quatre chantiers, par ordre de priorité décroissante : un
\textbf{second facteur d'authentification} sur le compte de direction, qui donne
accès à l'ensemble du fichier clientes ; un \textbf{test d'intrusion externe},
seule façon de découvrir ce que la personne qui a écrit le code ne pense pas à
chercher ; une \textbf{supervision} qui alerte sur les séries d'échecs de
connexion plutôt que de se contenter de les journaliser ; et le report du
compteur de débit dans un magasin partagé, préalable à toute mise à l'échelle
horizontale.

\paragraph{À plus long terme.}
\begin{itemize}[leftmargin=1.2cm,nosep]
  \item \textbf{Fiche cliente} tenue par l'institut : historique des soins,
  allergies signalées, préférences --- un dossier de suivi qui prolongerait la
  relation personnelle qui fait la valeur du centre.
  \item \textbf{Gestion multi-établissements}, si l'institut venait à ouvrir un
  second point de vente.
  \item \textbf{Automatisation du marketing} : relance des clientes non revenues
  depuis un certain temps, offres d'anniversaire.
\end{itemize}

\vspace{0.5cm}

Au terme de ce travail, la plate-forme est prête à entrer en service. Elle ne
remplace pas la relation entre l'institut et ses clientes --- elle ne le
cherche pas --- mais elle décharge la gérante de ce qui l'en éloignait : le
téléphone qui sonne pendant un soin, le carnet à consulter, la carte de
fidélité à retrouver. C'est, en définitive, la mesure la plus juste de sa
réussite.

% ═══════════════════════════════════════════════════════════════════
%  Webographie
% ═══════════════════════════════════════════════════════════════════
\pageLiminaire{Webographie}

Les ressources ci-dessous ont été consultées durant la réalisation du projet.
Les dates indiquées correspondent à la dernière consultation.

\vspace{0.4cm}

\begin{enumerate}[leftmargin=1.1cm,itemsep=6pt]

  \item \textbf{NestJS} --- documentation officielle du cadriciel de l'API
  (modules, gardes, intercepteurs, filtres d'exceptions, limitation de débit).
  \\ \url{https://docs.nestjs.com/}

  \item \textbf{Prisma} --- documentation de l'ORM : schéma, migrations,
  transactions, modes de suppression référentielle.
  \\ \url{https://www.prisma.io/docs}

  \item \textbf{PostgreSQL} --- documentation des fonctions de verrouillage
  consultatif, primitive au c\oe ur du moteur de réservation.
  \\ \url{https://www.postgresql.org/docs/current/explicit-locking.html}

  \item \textbf{React} --- documentation officielle de la bibliothèque
  d'interface du back-office.
  \\ \url{https://react.dev/}

  \item \textbf{React Native} --- documentation officielle du cadriciel de
  l'application mobile.
  \\ \url{https://reactnative.dev/docs/getting-started}

  \item \textbf{Expo} --- documentation de la chaîne d'outils mobile : stockage
  sécurisé, polices, images, compilation native.
  \\ \url{https://docs.expo.dev/}

  \item \textbf{TypeScript} --- manuel de référence du langage.
  \\ \url{https://www.typescriptlang.org/docs/}

  \item \textbf{Vite} --- documentation de l'outil de construction du
  back-office.
  \\ \url{https://vite.dev/guide/}

  \item \textbf{OWASP Top 10} --- classement des dix risques de sécurité
  applicative les plus critiques ; grille de référence du chapitre 4.
  \\ \url{https://owasp.org/www-project-top-ten/}

  \item \textbf{OWASP Cheat Sheet Series} --- \textit{Password Storage},
  \textit{Authentication}, \textit{Session Management}, \textit{File Upload} et
  \textit{Mass Assignment} : recommandations suivies pour le hachage des mots de
  passe, la défense contre l'énumération de comptes, la rotation des jetons et
  le contrôle des fichiers téléversés.
  \\ \url{https://cheatsheetseries.owasp.org/}

  \item \textbf{OWASP ASVS} --- \textit{Application Security Verification
  Standard} : référentiel de vérification employé pour établir la synthèse des
  contre-mesures.
  \\ \url{https://owasp.org/www-project-application-security-verification-standard/}

  \item \textbf{Microsoft --- Threat Modeling (STRIDE)} --- description du
  modèle de menaces appliqué au chapitre 4.
  \\ \url{https://learn.microsoft.com/en-us/azure/security/develop/threat-modeling-tool-threats}

  \item \textbf{MITRE CWE} --- \textit{Common Weakness Enumeration} :
  nomenclature des faiblesses logicielles citées (IDOR, injection, condition de
  course, exposition d'informations).
  \\ \url{https://cwe.mitre.org/}

  \item \textbf{JSON Web Tokens} --- introduction au format et à ses usages.
  \\ \url{https://jwt.io/introduction}

  \item \textbf{RFC 6819} --- \textit{OAuth 2.0 Threat Model and Security
  Considerations} : source de la stratégie de rotation des jetons de
  rafraîchissement et de la détection de réutilisation.
  \\ \url{https://datatracker.ietf.org/doc/html/rfc6819}

  \item \textbf{Argon2} --- spécification de la fonction de hachage retenue
  (RFC 9106).
  \\ \url{https://datatracker.ietf.org/doc/html/rfc9106}

  \item \textbf{RFC 2606} --- noms de domaine réservés, dont le domaine
  \texttt{.invalid} employé lors de l'anonymisation des comptes.
  \\ \url{https://datatracker.ietf.org/doc/html/rfc2606}

  \item \textbf{Helmet} --- documentation des en-têtes de sécurité HTTP posés
  par l'API.
  \\ \url{https://helmetjs.github.io/}

  \item \textbf{MDN --- HTTP security headers} --- référence sur HSTS, CSP,
  \texttt{X-Content-\allowbreak Type-Options}, CORS et CORP.
  \\ \url{https://developer.mozilla.org/en-US/docs/Web/HTTP/Headers}

  \item \textbf{IANA Time Zone Database} --- base des fuseaux horaires,
  référence de la zone \texttt{Africa/Casablanca}.
  \\ \url{https://www.iana.org/time-zones}

  \item \textbf{date-fns} --- documentation de la bibliothèque de manipulation
  de dates et de son extension pour les fuseaux horaires.
  \\ \url{https://date-fns.org/docs/Getting-Started}

  \item \textbf{OpenAPI / Swagger} --- spécification du format de documentation
  d'API et de son interface interactive.
  \\ \url{https://swagger.io/specification/}

  \item \textbf{Docker} --- documentation de la conteneurisation et de Docker
  Compose.
  \\ \url{https://docs.docker.com/}

  \item \textbf{GitHub Actions} --- documentation de la chaîne d'intégration
  continue.
  \\ \url{https://docs.github.com/en/actions}

  \item \textbf{Jest} --- documentation du cadriciel de tests.
  \\ \url{https://jestjs.io/docs/getting-started}

  \item \textbf{Scrum Guide} --- guide de référence de la méthode SCRUM.
  \\ \url{https://scrumguides.org/}

  \item \textbf{UML} --- spécification du langage de modélisation unifié
  publiée par l'OMG.
  \\ \url{https://www.omg.org/spec/UML/}

  \item \textbf{Lucidchart} --- plate-forme de création des diagrammes UML du
  présent rapport.
  \\ \url{https://www.lucidchart.com/}

  \item \textbf{CNDP} --- Commission Nationale de contrôle de la protection des
  Données à caractère Personnel : loi 09-08, obligations du responsable de
  traitement et droits des personnes concernées (accès, rectification,
  opposition, effacement).
  \\ \url{https://www.cndp.ma/}

  \item \textbf{i18next} --- documentation de la bibliothèque
  d'internationalisation employée côté mobile.
  \\ \url{https://www.i18next.com/}

\end{enumerate}


\end{document}
